%% ===== 自動生成：tools/mkguidetables.py（手で直さない。材料を直したら動かし直す） =====
%% ガイドブック（保護者・先生のための本）の付録（組 G-C、頭文字 kgc）。材料：tools/guide-taiou.md・series.py
\makeatletter
%% G-B の見出し（guide-p2a.tex の写し。p2a が先に読まれていれば、そちらが使われる）
%% \kgbsec[柱・しおり用の題]{番号}{題}（改ページしない）、\kgbsub{番号}{題}
\providecommand\kgbsec[3][]{%
  \par\addvspace{18pt}\Needspace{8\baselineskip}\global\bk@decotrue
  \setcounter{bkfig}{0}\setcounter{bktab}{0}\gdef\kgbsecno{#2}%
  \if\relax\detokenize{#1}\relax\def\kgb@st{#3}\else\def\kgb@st{#1}\fi
  \bksetrunning{}{保護者・先生のための本}{\if\relax\detokenize{#2}\relax\else#2\hspace{0.6\zw}\fi\kgb@st}%
  \pdfbookmark[1]{#2　\kgb@st}{kgbsec-#2-\thepage}%
  \kgc@secband{#2}{#3}}
\newcommand\kgc@secband[2]{%
  \noindent\begin{tikzpicture}
    \if\relax\detokenize{#1}\relax
      \node[anchor=west, text=white, font=\sffamily\bfseries\LARGE, align=left, inner sep=0pt] (t) at (0.6,0) {#2};
    \else
      \node[anchor=west, circle, fill=white, text=bkVol0, minimum size=1.35cm, font=\sffamily\bfseries\Large, inner sep=0pt] (n) at (0.35,0) {#1};
      \node[anchor=west, text=white, font=\sffamily\bfseries\LARGE, text width=\linewidth-2.4cm, align=left, inner sep=0pt] (t) at ([xshift=0.4cm]n.east) {#2};
    \fi
    \begin{pgfonlayer}{background}
      \path let \p1=(t.north), \p2=(t.south) in coordinate (kgbT) at (0,{max(\y1+0.4cm,1cm)}) coordinate (kgbB) at (0,{min(\y2-0.4cm,-1cm)});
      \fill[bkVol0] (kgbB) rectangle (\linewidth,0 |- kgbT);
    \end{pgfonlayer}
  \end{tikzpicture}\par\vspace{10pt}}
\providecommand\kgbsub[2]{%
  \par\addvspace{12pt}\Needspace{8\baselineskip}\noindent
  \begin{tikzpicture}[baseline=(t.base)]
    \node[anchor=west, fill=bkVol0, text=white, rounded corners=1mm, font=\sffamily\bfseries\small, inner xsep=1.6mm, inner ysep=1mm] (i) at (0,0) {#1};
    \node[anchor=west, font=\sffamily\bfseries\large, text=bkVol0, inner sep=0pt] (t) at ([xshift=0.6\zw]i.east) {#2};
    \draw[bkVol0, line width=0.8pt] ([yshift=-4pt]i.south west) -- ++(\linewidth-1pt,0);
  \end{tikzpicture}\par\nobreak\vspace{6pt}}
\providecommand\kgbsecno{}
%% ---- G-C の部品（kgc） ----
% 表の中の章番号（太いゴシック）
\newcommand\kgcch[1]{\mbox{\sffamily\bfseries #1}}
% 「もっと知りたい」（発展の囲み）の札：白い地に枠と字（色だけにたよらない）
\newcommand\kgcmotto{\mbox{\tikz[baseline=(t.base)]\node[draw=bkVol0, line width=0.5pt, rounded corners=0.8mm, inner xsep=0.5mm, inner ysep=0.35mm, font=\sffamily\footnotesize, text=bkVol0](t){もっと};}}
% 扱い方の記号の札（枠つき・太字）
\newcommand\kgcsym[1]{\mbox{\tikz[baseline=(t.base)]\node[draw=bkInk, line width=0.5pt, inner xsep=0.6mm, inner ysep=0.4mm, font=\sffamily\bfseries\footnotesize](t){#1};}}
\newcommand\kgcsep{\hspace{0.2em}}
\newcommand\kgcor{\hspace{0.15em}{\footnotesize ／}\hspace{0.15em}}
% ★：塗りの★＝その章の★の下から、白抜きの☆＝上まで（★★☆＝★★〜★★★）
\newcommand\kgcstars[2]{\mbox{\footnotesize\@tempcnta=0\relax
  \loop\advance\@tempcnta\@ne
    \ifnum\@tempcnta>#1\relax\textcolor{bkStar!70!black}{☆}\else\textcolor{bkStar}{★}\fi
  \ifnum\@tempcnta<#2\repeat}}
% 路線：アイコン（形の記号）＋丸数字
\newcommand\kgcroute[1]{\mbox{\bkrouteicon[0.95em]{#1}\,{\footnotesize\ifcase#1\or①\or②\or③\or④\or⑤\or⑥\or⑦\fi}}}
\newcommand\kgcallroutes{{\footnotesize ①〜⑦\newline すべて}}
% 範囲外の印（塗りの丸＝中心、半分塗りの丸＝一部。形で区別する）
\newcommand\kgcout[1]{\tikz[baseline=-0.35em]{\ifx#1a\fill[bkInk] (0,0) circle (0.42em);\else
  \fill[bkInk] (0,-0.42em) arc[start angle=-90, end angle=90, radius=0.42em] -- cycle;\fi
  \draw[bkInk, line width=0.6pt] (0,0) circle (0.42em);}}

% 凡例の囲み
\newenvironment{kgclegend}[1]{\par\addvspace{6pt}\begin{tcolorbox}[enhanced, breakable, arc=2mm, boxrule=0.8pt,
  colframe=bkVol0, colback=white, left=3mm, right=3mm, top=2.5mm, bottom=2.5mm,
  title={#1}, fonttitle=\sffamily\bfseries\small, coltitle=white, colbacktitle=bkVol0]\small}{\end{tcolorbox}\par\addvspace{4pt}}
% 表の見出しの行の色（白い地に濃い字の行より、うすい紺灰色の地に濃い字）
\colorlet{kgcHead}{bkVol0!14}
\colorlet{kgcBand}{bkVol0!6}
% longtblr の「表 1」などの見出しと続きの注を出さない
\DefTblrTemplate{caption}{kgc}{}
\DefTblrTemplate{capcont}{kgc}{\noindent{\sffamily\bfseries\small\color{bkVol0}\InsertTblrText{caption}}\hspace{0.6\zw}{\sffamily\footnotesize（前のページからの続き）}\par\vspace{2pt}}
\DefTblrTemplate{contfoot}{kgc}{\raggedleft\footnotesize（次のページに続きます）\par}
\NewTblrTheme{kgc}{\SetTblrTemplate{caption}{kgc}\SetTblrTemplate{capcont}{kgc}\SetTblrTemplate{contfoot}{kgc}}
\makeatother

\bkevenpage\kgbsec[付録A　学習指導要領の項目からシリーズの章へ]{A}{付録A　学習指導要領の項目から\\シリーズの章へ}\label{kgc-a}\label{kgc-sec10}
この付録は、現行の学習指導要領（小学校・中学校は平成29年告示、高等学校は平成30年告示）の算数・数学の項目ごとに、このシリーズのどの章で扱っているか、どのくらい扱っているかを示した一覧です。小学校第1学年から順に、中学校、高等学校の数学Ⅰ・数学A・数学Ⅱ・数学B・数学Ⅲ・数学Cの順に並べています。

\begin{itemize}[leftmargin=1.2\zw, itemsep=1pt, topsep=3pt]
\item 「番号」は、学年（科目）・領域・項目の番号です。例えば「小3-A(4)ア(ｴ)」は、小学校第3学年のA（数と計算）の(4)のアの(ｴ)を表します。「語」は〔用語・記号〕、「扱」は「内容の取扱い」の番号です。小学校のCは、第1〜3学年が「測定」、第4〜6学年が「変化と関係」です。
\item 「項目（要点）」は、指導要領の言葉を短くしたものです。正確な文は、指導要領の本文をご覧ください。
\item 「扱っている章」の \kgcch{1-4} などは章の番号（巻-章）です。「」の中は、その章の小見出しなどの題です。\kgcmotto{} は、章の中の発展の囲み「もっと知りたい」を表します。
\item 高等学校の科目を履修する学年は学校によって異なります。一般的な例は、数学Ⅰ・数学Aが1年、数学Ⅱ・数学Bが2年、数学Ⅲ・数学Cが3年です（指導要領で決められているものではありません）。
\item 「なし」「一部」の項目を補う方法は、「含めていない・軽く扱う内容と、補い方」の章をご覧ください。
\end{itemize}

\begin{kgclegend}{「扱い」の記号}
\begin{tblr}{colspec={Q[c,wd=10mm] X[l] Q[c,wd=10mm] X[l]}, rowsep=1.5pt, colsep=3pt, cells={font=\small}}
\kgcsym{同} & 指導要領とほぼ同じ範囲・深さで扱う & \kgcsym{後} & 後回し（その巻の「学年の目安」が、指導要領の学年より上。例：小3の内容を第4巻で） \\
\kgcsym{詳} & 詳しく（指導要領より深い・広い。証明を2通りにする、一般の形まで進む など） & \kgcsym{一部} & 一部だけ（語・考えの一部、「もっと知りたい」だけ、コラムだけ など） \\
\kgcsym{先} & 先取り（その巻の「学年の目安」が、指導要領の学年より下） & \kgcsym{なし} & 扱わない \\
\end{tblr}\par\smallskip
「先」「後」は、その巻の学年の目安と指導要領の学年を比べたもので、内容の深さ（同・詳・一部）とは別に書いています（例：\kgcsym{同}\kgcsep\kgcsym{後}）。「\kgcsym{一部}\kgcor\kgcsym{なし}」は、項目の中の事項によって扱いが分かれることを表します。\kgcsym{範囲外}は、指導要領の本文にない内容（数学Ⅲの注の微分方程式）を、このシリーズで扱っていることを表します。
\end{kgclegend}
\Needspace{259pt}\kgbsub{A.1}{小学校　第1学年}
\begin{longtblr}[theme=kgc, caption={小学校　第1学年}, entry=none, label=none]{colspec={Q[l,wd=19mm] Q[l,wd=31mm] X[2.4,l] Q[l,wd=15mm] X[1.6,l]}, width=\linewidth,
  rowhead=1, row{1}={bg=kgcHead, font=\sffamily\bfseries\footnotesize}, row{even[2]}={bg=kgcBand},
  hlines={0.4pt,bkLine}, vline{2-5}={0.4pt,bkLine}, rowsep=2pt, colsep=3pt, cells={font=\footnotesize}}
番号 & 項目（要点） & 扱っている章 & 扱い & 補足 \\
{\sffamily 小1-\allowbreak{}A(1)} & 数の構成と表し方 & \kgcch{1-1}（1対1対応・数える・順序と大小・「どれも3」）、\kgcch{1-2}（10のまとまり・一の位・十の位・1〜100の表・100・1000・10000）、\kgcch{1-6}（数直線）、\kgcch{1-4}（いくつといくつ＝ア(ｴ)） & \kgcsym{同} & \kgcch{1-2} は 1000・10000 まで進む（小2〜小3の範囲も含む） \\
{\sffamily 小1-\allowbreak{}A(2)} & 加法及び減法 & \kgcch{1-4}（あわせる・ふえる・10をこえる足し算・0の足し算）、\kgcch{1-5}（のこり・ちがい・かわった量・くり下がり）、\kgcch{1-7}（何十の計算） & \kgcsym{同} & \kgcch{1-5} は「かわった量」（変化量）を足している（路線①の始発） \\
{\sffamily 小1-\allowbreak{}B(1)} & 図形（形） & \kgcch{1-19}（三角形・四角形・円、つつ・はこを転がす） & \kgcsym{一部} & ア(ｳ)の方向や位置を表す言葉（前後・左右・上下）を教える節はなく、右・左は \kgcch{1-6} の数直線（右へ進む・左へ戻る）や \kgcch{1-2} の位（左の位）で使うだけ \\
{\sffamily 小1-\allowbreak{}C(1)} & 量の大きさの比較 & \kgcch{1-1}（ペアにして比べる）、\kgcch{1-8}（1 cm のいくつ分）、\kgcch{1-9}（てんびんで重さを比べる）、\kgcch{1-21}（マスの数で広さを比べる） & \kgcsym{同} &  \\
{\sffamily 小1-\allowbreak{}C(2)} & 時刻の読み方 & \kgcch{1-10} & \kgcsym{同} &  \\
{\sffamily 小1-\allowbreak{}D(1)} & 数量の整理（絵や図） & \kgcch{1-17}（数えて表に・ぼうグラフ） & \kgcsym{同} & 「絵グラフ」の語はない \\
{\sffamily 小1-\allowbreak{}語} & 一の位、十の位、＋、－、＝ & \kgcch{1-2}・\kgcch{1-4}・\kgcch{1-5} & \kgcsym{同} & 記号カード：＋・＝（\kgcch{1-4}）、－（\kgcch{1-5}） \\
\end{longtblr}
\Needspace{229pt}\kgbsub{A.2}{小学校　第2学年}
\begin{longtblr}[theme=kgc, caption={小学校　第2学年}, entry=none, label=none]{colspec={Q[l,wd=19mm] Q[l,wd=31mm] X[2.4,l] Q[l,wd=15mm] X[1.6,l]}, width=\linewidth,
  rowhead=1, row{1}={bg=kgcHead, font=\sffamily\bfseries\footnotesize}, row{even[2]}={bg=kgcBand},
  hlines={0.4pt,bkLine}, vline{2-5}={0.4pt,bkLine}, rowsep=2pt, colsep=3pt, cells={font=\footnotesize}}
番号 & 項目（要点） & 扱っている章 & 扱い & 補足 \\
{\sffamily 小2-\allowbreak{}A(1)} & 数の構成と表し方（4位数まで・十や百を単位・積としてみる・簡単な分数） & \kgcch{1-2}（1000・10000）、\kgcch{1-6}（1目もり10・100の数直線）、\kgcch{1-12}・\kgcch{1-13}（積としてみる）、\kgcch{1-16}（1/2・1/4） & \kgcsym{同} &  \\
{\sffamily 小2-\allowbreak{}A(2)} & 加法及び減法（2位数・筆算・性質・相互関係） & \kgcch{1-7}「ひっ算と くり上がり」「ひっ算と くり下がり」「こたえを たしかめよう」 & \kgcsym{同} &  \\
{\sffamily 小2-\allowbreak{}A(3)} & 乗法（意味・九九・性質・2位数×1位数） & \kgcch{1-12}（意味・アレイ・2〜5のだん）、\kgcch{1-13}（九九の表・ふえ方・入れかえても同じ・だんをたしてつくる・3×10） & \kgcsym{同}\newline{}（ア(ｵ)は一部） & ア(ｵ)「簡単な場合の2位数と1位数の乗法」は、\kgcch{1-13} の「3×10」、\kgcch{1-20} の「30×4」程度 \\
{\sffamily 小2-\allowbreak{}B(1)} & 三角形や四角形などの図形（箱の形を含む） & \kgcch{1-19}「三角形と 四角形」「直角」「直角の ある 形」 & \kgcsym{同}\newline{}（ア(ｳ)は一部） & ア(ｳ)「箱の形」は、\kgcch{1-19} のはてな？で「はこ」を転がす程度 \\
{\sffamily 小2-\allowbreak{}C(1)} & 長さ（mm・cm・m）、かさ（mL・dL・L） & \kgcch{1-8}、\kgcch{1-9} & \kgcsym{同} &  \\
{\sffamily 小2-\allowbreak{}C(2)} & 時間の単位（日・時・分） & \kgcch{1-10}「60分で 1時間」「午前と 午後」 & \kgcsym{同} &  \\
{\sffamily 小2-\allowbreak{}D(1)} & 簡単な表やグラフ & \kgcch{1-17} & \kgcsym{同} &  \\
{\sffamily 小2-\allowbreak{}語} & 直線、直角、頂点、辺、面、単位、×、＞、＜ & \kgcch{1-19}（辺・ちょう点・直角）、\kgcch{1-12}（×）、\kgcch{2-2}（＜・＞） & \kgcsym{一部}\kgcsep{}\kgcsym{後} & ＞・＜は第1巻にない（第2巻 \kgcch{2-2} で不等号） \\
{\sffamily 小2-\allowbreak{}扱} & (2)（ ）や□の式 (3)交換・結合法則 (4)乗数が1ずつ増える (5)敷き詰め & (3)(4)＝\kgcch{1-13}「入れかえても同じ」「ふえかたの きまり」、\kgcch{2-11}（結合・分配）。(2)＝\kgcch{2-11} & \kgcsym{一部} &  \\
\end{longtblr}
\Needspace{184pt}\kgbsub{A.3}{小学校　第3学年}
\begin{longtblr}[theme=kgc, caption={小学校　第3学年}, entry=none, label=none]{colspec={Q[l,wd=19mm] Q[l,wd=31mm] X[2.4,l] Q[l,wd=15mm] X[1.6,l]}, width=\linewidth,
  rowhead=1, row{1}={bg=kgcHead, font=\sffamily\bfseries\footnotesize}, row{even[2]}={bg=kgcBand},
  hlines={0.4pt,bkLine}, vline{2-5}={0.4pt,bkLine}, rowsep=2pt, colsep=3pt, cells={font=\footnotesize}}
番号 & 項目（要点） & 扱っている章 & 扱い & 補足 \\
{\sffamily 小3-\allowbreak{}A(1)} & 整数の表し方（万の単位、10倍・1/10） & \kgcch{1-2}「1000と 10000」、\kgcch{2-1}「10倍すると、小数点が右へ1つ」、\kgcch{2-7}（万・億・兆の表） & \kgcsym{一部}\kgcsep{}\kgcsym{後} & 「万の単位」の読み書きの節は第1巻にない。万→\nobreak{}億→\nobreak{}兆は \kgcch{2-7} \\
{\sffamily 小3-\allowbreak{}A(2)} & 加法及び減法（3位数・4位数） & \kgcch{1-7}「100を こえる たし算」 & \kgcsym{同} &  \\
{\sffamily 小3-\allowbreak{}A(3)} & 乗法（2位数・3位数×1位数・2位数、筆算） & \kgcch{1-13}（3×10）、\kgcch{1-20}（30×2〜30×6）。筆算の形は \kgcch{3-8} のコラム「かけ算の筆算も、展開だった」（23×14）だけ & \kgcsym{一部} & \textbf{掛け算の筆算の節はない} \\
{\sffamily 小3-\allowbreak{}A(4)} & 除法（意味・余り・乗法との関係・九九で求める・商が2位数） & \kgcch{1-14}（等分除・包含除・九九で見つける・1でわる・0では分けられない）、\kgcch{1-15}（余り・確かめの式） & \kgcsym{同}\newline{}（ア(ｵ)はなし） & 除法の例は九九の範囲（10÷2〜32÷8 など）。ア(ｵ)「商が2位数」はない \\
{\sffamily 小3-\allowbreak{}A(5)} & 小数（1/10 の位、加減） & \kgcch{1-16}「1を 10に 分けた 1つ分：0.1」、\kgcch{2-1}（1/10・1/100 の位） & \kgcsym{同} & 計算は \kgcch{2-4} \\
{\sffamily 小3-\allowbreak{}A(6)} & 分数（等分・単位分数・簡単な加減） & \kgcch{1-16}（1/2・1/4、0.5と1/2は同じ）、\kgcch{2-4}（足し算・引き算） & \kgcsym{一部}\kgcsep{}\kgcsym{後} & 第1巻は 1/2・1/4 まで \\
{\sffamily 小3-\allowbreak{}A(7)} & 数量の関係を表す式（□） & \kgcch{2-11}「ことばの式を、□で短く書く」 & \kgcsym{同}\kgcsep{}\kgcsym{後} &  \\
{\sffamily 小3-\allowbreak{}A(8)} & そろばん & なし & \kgcsym{なし} & どの巻にも語がない \\
{\sffamily 小3-\allowbreak{}B(1)} & 二等辺三角形・正三角形、作図、角、円・球 & 円・中心・半径・直径・コンパス＝\kgcch{1-19}「円と タイヤ」。角＝\kgcch{1-20}。二等辺三角形・正三角形＝\kgcch{4-1}（対称で性質を説明）。球＝\kgcch{4-5} & \kgcsym{一部}\kgcsep{}\kgcsym{後} & 定規とコンパスで二等辺三角形をかく作図はない \\
{\sffamily 小3-\allowbreak{}C(1)} & 長さ（km）、重さ（g・kg）、測定 & \kgcch{1-9}「おもさの たんい：g と kg」、\kgcch{1-11}（目もり・はかり）、km は \kgcch{2-10}（接頭語 k） & \kgcsym{一部}\kgcsep{}\kgcsym{後} & 第1巻に km はない \\
{\sffamily 小3-\allowbreak{}C(2)} & 時刻と時間（秒、時間を求める） & \kgcch{1-10}「時こくと 時間」「\kgcmotto{} みじかい 時間：秒」 & \kgcsym{同} &  \\
{\sffamily 小3-\allowbreak{}D(1)} & 表と棒グラフ & \kgcch{1-17}「1目もりの 大きさ」「1目もりの きめかた」 & \kgcsym{同} &  \\
{\sffamily 小3-\allowbreak{}語} & 等号、不等号、小数点、1/10の位、数直線、分母、分子、÷ & \kgcch{1-4}（＝）、\kgcch{1-6}（数直線）、\kgcch{1-14}（÷）、\kgcch{1-16}（小数点）、\kgcch{2-1}（1/10の位）、\kgcch{2-2}（不等号）、\kgcch{2-4}（分母・分子） & \kgcsym{同}\newline{}（不等号は後） &  \\
{\sffamily 小3-\allowbreak{}扱} & (7)トン（t）、接頭語（k・m） & \kgcch{2-10}「大きさを表す記号：k・m・$\symup{\mu}$」 & \kgcsym{詳}\kgcsep{}\kgcsym{後} &  \\
\end{longtblr}
\Needspace{154pt}\kgbsub{A.4}{小学校　第4学年}
\begin{longtblr}[theme=kgc, caption={小学校　第4学年}, entry=none, label=none]{colspec={Q[l,wd=19mm] Q[l,wd=31mm] X[2.4,l] Q[l,wd=15mm] X[1.6,l]}, width=\linewidth,
  rowhead=1, row{1}={bg=kgcHead, font=\sffamily\bfseries\footnotesize}, row{even[2]}={bg=kgcBand},
  hlines={0.4pt,bkLine}, vline{2-5}={0.4pt,bkLine}, rowsep=2pt, colsep=3pt, cells={font=\footnotesize}}
番号 & 項目（要点） & 扱っている章 & 扱い & 補足 \\
{\sffamily 小4-\allowbreak{}A(1)} & 整数の表し方（億・兆） & \kgcch{2-7}「大きな数は、10の累乗で」 & \kgcsym{同} &  \\
{\sffamily 小4-\allowbreak{}A(2)} & 概数（四捨五入・見積り） & \kgcch{2-6}「がい数と四捨五入」 & \kgcsym{同} & 「概数」の語はなく「がい数」 \\
{\sffamily 小4-\allowbreak{}A(3)} & 整数の除法（筆算、被除数＝除数×商＋余り） & ア(ｳ)＝\kgcch{1-15}「たしかめの しき」 & \kgcsym{一部} & \textbf{割り算の筆算の節はない} \\
{\sffamily 小4-\allowbreak{}A(4)} & 小数の仕組みと計算 & \kgcch{2-1}、\kgcch{2-4} & \kgcsym{同} &  \\
{\sffamily 小4-\allowbreak{}A(5)} & 分数（等しい分数・同分母の加減） & \kgcch{2-4}「大きさの同じ分数：約分と通分」 & \kgcsym{同} & 真分数・仮分数・帯分数の語も \kgcch{2-4} \\
{\sffamily 小4-\allowbreak{}A(6)} & 数量の関係を表す式（四則混合・（ ）・公式・□△） & \kgcch{2-11}、\kgcch{2-13} & \kgcsym{同} &  \\
{\sffamily 小4-\allowbreak{}A(7)} & 計算に関して成り立つ性質 & \kgcch{2-11}「文字で、いつでも成り立つきまりを書く」、\kgcch{1-13} & \kgcsym{同} &  \\
{\sffamily 小4-\allowbreak{}A(8)} & そろばん & なし & \kgcsym{なし} &  \\
{\sffamily 小4-\allowbreak{}B(1)} & 平面図形（平行・垂直、平行四辺形・ひし形・台形） & \kgcch{4-1}（⊥・∥、平行線と角）、\kgcch{4-5}（平行四辺形・台形・ひし形の面積） & \kgcsym{同}\kgcsep{}\kgcsym{後} &  \\
{\sffamily 小4-\allowbreak{}B(2)} & 立体図形（立方体・直方体、見取図・展開図） & \kgcch{2-7}（見取図）、\kgcch{4-5}「展開図と表面積・回転体」、\kgcch{4-13}（投影図・空間の平行垂直） & \kgcsym{同}\kgcsep{}\kgcsym{後} &  \\
{\sffamily 小4-\allowbreak{}B(3)} & ものの位置の表し方 & \kgcch{3-2}（座標）、\kgcch{3-2} \kgcmotto{}「部屋の中を飛ぶハエは、3つの数」、\kgcch{4-13} & \kgcsym{詳}\kgcsep{}\kgcsym{後} &  \\
{\sffamily 小4-\allowbreak{}B(4)} & 平面図形の面積（cm$^{2}$・m$^{2}$・km$^{2}$、正方形・長方形） & \kgcch{1-21}「1cm$^{2}$」「長方形の 面積は かけ算で」「大きな 広さ：1m$^{2}$」、\kgcch{2-10}「面積・体積の換算」 & \kgcsym{同}\kgcsep{}\kgcsym{先} & 小4-扱(8) アール・ヘクタールはない \\
{\sffamily 小4-\allowbreak{}B(5)} & 角の大きさ（度・測定） & \kgcch{1-20} & \kgcsym{同}\kgcsep{}\kgcsym{先} &  \\
{\sffamily 小4-\allowbreak{}C(1)} & 伴って変わる二つの数量（折れ線グラフ） & \kgcch{1-18} & \kgcsym{同}\kgcsep{}\kgcsym{先} & 「かたむきが 急なほど、はやく かわる」は傾き（中2）の種 \\
{\sffamily 小4-\allowbreak{}C(2)} & 簡単な場合についての割合 & \kgcch{2-5}「割合：もとにする量を1とみる」 & \kgcsym{同} &  \\
{\sffamily 小4-\allowbreak{}D(1)} & データの分類整理（二つの観点・折れ線グラフ） & ア(ｲ)＝\kgcch{1-18} & \kgcsym{一部} &  \\
{\sffamily 小4-\allowbreak{}語} & 和・差・積・商、以上・以下・未満、真分数・仮分数・帯分数、平行・垂直・対角線・平面 & \kgcch{2-2}（以上・以下・未満）、\kgcch{2-4}（真分数など）、\kgcch{4-1}（平行・垂直）、\kgcch{4-13}（平面） & \kgcsym{同} &  \\
\end{longtblr}
\Needspace{259pt}\kgbsub{A.5}{小学校　第5学年}
\begin{longtblr}[theme=kgc, caption={小学校　第5学年}, entry=none, label=none]{colspec={Q[l,wd=19mm] Q[l,wd=31mm] X[2.4,l] Q[l,wd=15mm] X[1.6,l]}, width=\linewidth,
  rowhead=1, row{1}={bg=kgcHead, font=\sffamily\bfseries\footnotesize}, row{even[2]}={bg=kgcBand},
  hlines={0.4pt,bkLine}, vline{2-5}={0.4pt,bkLine}, rowsep=2pt, colsep=3pt, cells={font=\footnotesize}}
番号 & 項目（要点） & 扱っている章 & 扱い & 補足 \\
{\sffamily 小5-\allowbreak{}A(1)} & 整数の性質（偶数・奇数、約数・倍数） & \kgcch{2-4}（公約数・公倍数・素数）、\kgcch{2-9} \kgcmotto{}「$\sqrt{2}$ が分数で書けないわけ」（偶数・奇数の定義） & \kgcsym{一部} & 偶数・奇数を定義するのは \kgcch{2-9} の「もっと知りたい」（\mbox{★★★}）だけ。「最大公約数」「最小公倍数」の語はない（「公倍数でいちばん小さい数」と言う） \\
{\sffamily 小5-\allowbreak{}A(2)} & 整数及び小数の表し方（小数点の移動） & \kgcch{2-1} & \kgcsym{同} &  \\
{\sffamily 小5-\allowbreak{}A(3)} & 小数の乗法及び除法 & \kgcch{2-4}「0と1のあいだの数をかける・わる」 & \kgcsym{同} &  \\
{\sffamily 小5-\allowbreak{}A(4)} & 分数の意味と表し方 & \kgcch{2-4}「分数は、わり算の答え」 & \kgcsym{同} &  \\
{\sffamily 小5-\allowbreak{}A(5)} & 分数の加法及び減法（異分母） & \kgcch{2-4}「たし算・ひき算は、通分してから」 & \kgcsym{同} &  \\
{\sffamily 小5-\allowbreak{}A(6)} & 数量の関係を表す式 & \kgcch{2-11}、\kgcch{2-15} & \kgcsym{同} &  \\
{\sffamily 小5-\allowbreak{}B(1)} & 平面図形（合同・多角形・正多角形・円周率） & \kgcch{4-1}（合同と対称・多角形の内角と外角・正多角形）、\kgcch{2-9}「円周率 $\pi$」 & \kgcsym{同}\kgcsep{}\kgcsym{後}\newline{}（$\pi$ は同） & プログラミングで正多角形をかく例（小5-プログラミング）はない \\
{\sffamily 小5-\allowbreak{}B(2)} & 立体図形（角柱・円柱） & \kgcch{4-5}「柱の体積：底面積×高さ」 & \kgcsym{同}\kgcsep{}\kgcsym{後} &  \\
{\sffamily 小5-\allowbreak{}B(3)} & 平面図形の面積（三角形・平行四辺形・ひし形・台形） & \kgcch{4-5} & \kgcsym{同}\kgcsep{}\kgcsym{後} &  \\
{\sffamily 小5-\allowbreak{}B(4)} & 立体図形の体積（cm$^{3}$・m$^{3}$、直方体・立方体） & \kgcch{2-7}「面積は2乗、体積は3乗」、\kgcch{2-10}、\kgcch{4-5} & \kgcsym{同} &  \\
{\sffamily 小5-\allowbreak{}C(1)} & 比例（簡単な場合） & \kgcch{2-15} & \kgcsym{同} &  \\
{\sffamily 小5-\allowbreak{}C(2)} & 速さなど単位量当たりの大きさ & \kgcch{2-3}「わり算で『1あたりの量』」、\kgcch{2-10}「時速を秒速に直す」 & \kgcsym{同} &  \\
{\sffamily 小5-\allowbreak{}C(3)} & 割合・百分率 & \kgcch{2-5}「百分率：もとを100とみる」、\kgcch{2-5} で歩合 & \kgcsym{同} &  \\
{\sffamily 小5-\allowbreak{}D(1)} & 円グラフ・帯グラフ、統計的な問題解決 & なし & \kgcsym{なし} & 円グラフ・帯グラフはどの巻にもない。統計的な問題解決の手順（問い→\nobreak{}計画→\nobreak{}データ→\nobreak{}分析→\nobreak{}結論）としての節はない。データの集め方・表・グラフは \kgcch{1-17}・\kgcch{1-18}・\kgcch{3-14}〜\kgcch{3-16} などで扱う \\
{\sffamily 小5-\allowbreak{}D(2)} & 平均 & \kgcch{2-6}「平均：でこぼこをならして1つの数にする」、\kgcch{3-13} & \kgcsym{同} &  \\
{\sffamily 小5-\allowbreak{}語} & 最大公約数、最小公倍数、通分、約分、底面、側面、比例、％ & \kgcch{2-4}（通分・約分）、\kgcch{2-5}（％）、\kgcch{2-15}（比例）、\kgcch{4-5}（底面） & \kgcsym{一部} & 最大公約数・最小公倍数の語はない \\
\end{longtblr}
\Needspace{184pt}\kgbsub{A.6}{小学校　第6学年}
\begin{longtblr}[theme=kgc, caption={小学校　第6学年}, entry=none, label=none]{colspec={Q[l,wd=19mm] Q[l,wd=31mm] X[2.4,l] Q[l,wd=15mm] X[1.6,l]}, width=\linewidth,
  rowhead=1, row{1}={bg=kgcHead, font=\sffamily\bfseries\footnotesize}, row{even[2]}={bg=kgcBand},
  hlines={0.4pt,bkLine}, vline{2-5}={0.4pt,bkLine}, rowsep=2pt, colsep=3pt, cells={font=\footnotesize}}
番号 & 項目（要点） & 扱っている章 & 扱い & 補足 \\
{\sffamily 小6-\allowbreak{}A(1)} & 分数の乗法及び除法 & \kgcch{2-4}「かけ算は、面積図で」「わり算（その2）：なぜ、上と下を入れかえた数をかけるの？」 & \kgcsym{同} &  \\
{\sffamily 小6-\allowbreak{}A(2)} & 文字を用いた式 & \kgcch{2-11}「□のかわりに、文字を使う」 & \kgcsym{同} &  \\
{\sffamily 小6-\allowbreak{}B(1)} & 縮図・拡大図、対称な図形 & \kgcch{4-3}「拡大図・縮図と縮尺」、\kgcch{4-1}「ずらす・回す・裏返す：合同と対称」 & \kgcsym{同}\kgcsep{}\kgcsym{後} &  \\
{\sffamily 小6-\allowbreak{}B(2)} & 概形とおよその面積 & \kgcch{4-5}「方眼で数えて、細かくするほど正確に」、\kgcch{5-9} \kgcmotto{}「グラフが曲がっていても、マスを数えればおよそわかる」 & \kgcsym{同}\kgcsep{}\kgcsym{後} &  \\
{\sffamily 小6-\allowbreak{}B(3)} & 円の面積 & \kgcch{4-5}「円の面積：細いおうぎ形に切って並べかえる」 & \kgcsym{同}\kgcsep{}\kgcsym{後} & \kgcch{4-2} は「\kgcch{4-5} でくわしく」と先に使う \\
{\sffamily 小6-\allowbreak{}B(4)} & 角柱及び円柱の体積 & \kgcch{4-5} & \kgcsym{同}\kgcsep{}\kgcsym{後} &  \\
{\sffamily 小6-\allowbreak{}C(1)} & 比例（反比例は「知る」程度） & \kgcch{2-15}「かけると一定：反比例」、\kgcch{3-7}（反比例のグラフ） & \kgcsym{詳} & 反比例を xy 一定・グラフ・pV 一定まで扱う \\
{\sffamily 小6-\allowbreak{}C(2)} & 比 & \kgcch{2-5}「比と比の値」「比はかけ算でつなげる」 & \kgcsym{同} &  \\
{\sffamily 小6-\allowbreak{}D(1)} & データの考察（代表値・度数分布） & \kgcch{3-13}「1つの数で代表させる：平均・中央値・最頻値」「数を区切って数える：度数分布表」、\kgcch{2-6}（点を並べた図） & \kgcsym{同}\kgcsep{}\kgcsym{後} & \textbf{「ドットプロット」の語は紙面にない}（\kgcch{2-6} は点を並べた図で、名前は付けていない） \\
{\sffamily 小6-\allowbreak{}D(2)} & 起こり得る場合 & \kgcch{3-17}「場合を数えて、確率を求める」（表・樹形図） & \kgcsym{同}\kgcsep{}\kgcsym{後} &  \\
{\sffamily 小6-\allowbreak{}語} & 線対称・点対称・対称の軸・対称の中心、比の値、ドットプロット、平均値・中央値・最頻値・階級、： & \kgcch{4-1}（線対称・点対称）、\kgcch{2-5}（比の値・：）、\kgcch{3-13}（中央値・最頻値・階級） & \kgcsym{一部} & ドットプロットの語なし（上） \\
\end{longtblr}
\Needspace{169pt}\kgbsub{A.7}{中学校　第1学年}
\begin{longtblr}[theme=kgc, caption={中学校　第1学年}, entry=none, label=none]{colspec={Q[l,wd=19mm] Q[l,wd=31mm] X[2.4,l] Q[l,wd=15mm] X[1.6,l]}, width=\linewidth,
  rowhead=1, row{1}={bg=kgcHead, font=\sffamily\bfseries\footnotesize}, row{even[2]}={bg=kgcBand},
  hlines={0.4pt,bkLine}, vline{2-5}={0.4pt,bkLine}, rowsep=2pt, colsep=3pt, cells={font=\footnotesize}}
番号 & 項目（要点） & 扱っている章 & 扱い & 補足 \\
{\sffamily 中1-\allowbreak{}A(1)} & 正の数・負の数（扱(1)素数の積） & \kgcch{2-2}（正負の数の四則）、\kgcch{2-4}（素数の掛け算に分ける） & \kgcsym{同} & 「素因数分解」の語はない \\
{\sffamily 中1-\allowbreak{}A(2)} & 文字を用いた式 & \kgcch{2-11}、\kgcch{2-12}、\kgcch{2-13} & \kgcsym{同} &  \\
{\sffamily 中1-\allowbreak{}A(3)} & 一元一次方程式（扱(4)比例式） & \kgcch{2-14}「等式は天びん」「移項」、\kgcch{2-16}「1次方程式の解き方」、比例式＝\kgcch{3-3}「いろいろな関数」の中の段 & \kgcsym{同} &  \\
{\sffamily 中1-\allowbreak{}B(1)} & 平面図形（基本的な作図、移動） & 移動＝\kgcch{4-1}「ずらす・回す・裏返す」、\kgcch{4-10}「平行移動」。作図＝\kgcch{4-2}「垂直二等分線と、円の中心」（紙を折って見つける） & \kgcsym{一部}\kgcsep{}\kgcsym{後} & 角の二等分線・垂線の作図はない \\
{\sffamily 中1-\allowbreak{}B(2)} & 空間図形（位置関係、扇形、表面積と体積） & \kgcch{4-13}（ねじれの位置・投影図）、\kgcch{4-2}「おうぎ形の弧と面積は、中心角に比例する」、\kgcch{4-5}「すいの体積」「球の体積と表面積」 & \kgcsym{同}\kgcsep{}\kgcsym{後} &  \\
{\sffamily 中1-\allowbreak{}C(1)} & 比例・反比例（関数・座標・表式グラフ） & \kgcch{2-15}、\kgcch{3-1}（変数・変域）、\kgcch{3-2}（座標）、\kgcch{3-3}（関数）、\kgcch{3-7} & \kgcsym{同}\kgcsep{}\kgcsym{詳} &  \\
{\sffamily 中1-\allowbreak{}D(1)} & データの分布（ヒストグラム・相対度数） & \kgcch{3-13}「柱のグラフにする：ヒストグラム」 & \kgcsym{同}\kgcsep{}\kgcsym{後} & 累積度数も \kgcch{3-13} \\
{\sffamily 中1-\allowbreak{}D(2)} & 多数回の試行による確率 & \kgcch{3-17}「たくさんくり返すと、割合が落ち着く」 & \kgcsym{同}\kgcsep{}\kgcsym{後} & 表計算の RANDBETWEEN で試す \\
{\sffamily 中1-\allowbreak{}語} & 自然数・素数・符号・絶対値・項・係数・移項・≦≧、弧・弦・回転体・ねじれの位置・$\pi$・∥・⊥・∠・△、関数・変数・変域、範囲・累積度数 & \kgcch{2-2}・\kgcch{2-4}・\kgcch{2-9}・\kgcch{2-12}・\kgcch{2-14}・\kgcch{3-1}・\kgcch{3-2}・\kgcch{3-13}・\kgcch{4-1}・\kgcch{4-2}・\kgcch{4-5}・\kgcch{4-13} & \kgcsym{同} &  \\
{\sffamily 中1-\allowbreak{}扱} & (5)円の接線は半径に垂直 (6)投影図 & \kgcch{4-2}「円の接線は、半径に垂直」、\kgcch{4-13}「空間の形を紙にかく：投影図」 & \kgcsym{同}\kgcsep{}\kgcsym{後} &  \\
\end{longtblr}
\Needspace{184pt}\kgbsub{A.8}{中学校　第2学年}
\begin{longtblr}[theme=kgc, caption={中学校　第2学年}, entry=none, label=none]{colspec={Q[l,wd=19mm] Q[l,wd=31mm] X[2.4,l] Q[l,wd=15mm] X[1.6,l]}, width=\linewidth,
  rowhead=1, row{1}={bg=kgcHead, font=\sffamily\bfseries\footnotesize}, row{even[2]}={bg=kgcBand},
  hlines={0.4pt,bkLine}, vline{2-5}={0.4pt,bkLine}, rowsep=2pt, colsep=3pt, cells={font=\footnotesize}}
番号 & 項目（要点） & 扱っている章 & 扱い & 補足 \\
{\sffamily 中2-\allowbreak{}A(1)} & 文字を用いた式（整式の加減・単項式の乗除・式の変形） & \kgcch{2-12}、\kgcch{2-14}「x＝vt から、時間を求める式を作る」、\kgcch{3-8} & \kgcsym{同}\kgcsep{}\kgcsym{先} &  \\
{\sffamily 中2-\allowbreak{}A(2)} & 連立二元一次方程式 & \kgcch{3-6} & \kgcsym{同} & \kgcmotto{}「3つの文字の連立方程式」（\mbox{★★★}） \\
{\sffamily 中2-\allowbreak{}B(1)} & 平行線や角、多角形の角 & \kgcch{4-1}「平行線と角：対頂角・同位角・錯角」 & \kgcsym{同} &  \\
{\sffamily 中2-\allowbreak{}B(2)} & 図形の合同（合同条件・証明） & \kgcch{4-1}「三角形が1つに決まる条件：合同条件」「三角形の内角の和は180°」 & \kgcsym{一部} & 証明は「内角の和」の1題。二等辺三角形・平行四辺形の性質は対称で説明し、証明しない \\
{\sffamily 中2-\allowbreak{}C(1)} & 一次関数（二元一次方程式を関数とみる） & \kgcch{3-5}、\kgcch{3-6}「交わる点は、2つの式を同時にみたす」 & \kgcsym{同} &  \\
{\sffamily 中2-\allowbreak{}D(1)} & 四分位範囲・箱ひげ図 & \kgcch{3-13}「4つに分けて見る：四分位数と箱ひげ図」 & \kgcsym{同} &  \\
{\sffamily 中2-\allowbreak{}D(2)} & 場合の数を基にして得られる確率 & \kgcch{3-17} & \kgcsym{同} &  \\
{\sffamily 中2-\allowbreak{}語} & 同類項、対頂角・内角・外角・定義・証明・逆・反例・≡、変化の割合・傾き & \kgcch{2-12}、\kgcch{4-1}「逆と反例」、\kgcch{4-3}（≡ の記号カード）、\kgcch{3-5} & \kgcsym{同} &  \\
{\sffamily 中2-\allowbreak{}扱} & (1)正方形・ひし形・長方形は平行四辺形の特別な形 & \kgcch{4-1}（「ひし形は、正方形ではない」：逆と反例）、\kgcch{4-5}「台形・ひし形・多角形」（ひし形の定義と面積） & \kgcsym{一部} &  \\
\end{longtblr}
\Needspace{184pt}\kgbsub{A.9}{中学校　第3学年}
\begin{longtblr}[theme=kgc, caption={中学校　第3学年}, entry=none, label=none]{colspec={Q[l,wd=19mm] Q[l,wd=31mm] X[2.4,l] Q[l,wd=15mm] X[1.6,l]}, width=\linewidth,
  rowhead=1, row{1}={bg=kgcHead, font=\sffamily\bfseries\footnotesize}, row{even[2]}={bg=kgcBand},
  hlines={0.4pt,bkLine}, vline{2-5}={0.4pt,bkLine}, rowsep=2pt, colsep=3pt, cells={font=\footnotesize}}
番号 & 項目（要点） & 扱っている章 & 扱い & 補足 \\
{\sffamily 中3-\allowbreak{}A(1)} & 平方根 & \kgcch{2-8}（√・√どうしの掛け算）、\kgcch{2-9}（無理数）、\kgcch{3-10}「√ のついた数の計算」（有理化） & \kgcsym{同}\kgcsep{}\kgcsym{先} & \kgcch{2-8}・\kgcch{2-9} は\mbox{★★}の章（第2巻） \\
{\sffamily 中3-\allowbreak{}A(2)} & 式の展開と因数分解 & \kgcch{3-8}、\kgcch{3-10}「かけ算の形に直す：因数分解」 & \kgcsym{同} &  \\
{\sffamily 中3-\allowbreak{}A(3)} & 二次方程式（解の公式） & \kgcch{3-10} & \kgcsym{詳} & 判別式・2次不等式まで（数Ⅰ・数Ⅱの内容） \\
{\sffamily 中3-\allowbreak{}B(1)} & 図形の相似（相似条件・面積比・体積比、平行線と線分の比） & \kgcch{4-3} & \kgcsym{同} & 相似の証明の場面はない \\
{\sffamily 中3-\allowbreak{}B(2)} & 円周角と中心角 & \kgcch{4-2}「円周角は、中心角の半分」 & \kgcsym{同} &  \\
{\sffamily 中3-\allowbreak{}B(3)} & 三平方の定理 & \kgcch{4-4}「どんな直角三角形でも成り立つわけ：2通りの証明」 & \kgcsym{詳} & 指導要領は「証明できることを知る」。シリーズは2通りの証明を本文で見せる \\
{\sffamily 中3-\allowbreak{}C(1)} & 関数 y＝ax$^{2}$ & \kgcch{3-9} & \kgcsym{同} &  \\
{\sffamily 中3-\allowbreak{}D(1)} & 標本調査 & \kgcch{3-17}「全部を調べずに、一部から全体を知る：標本調査」（母集団・標本・無作為・全数調査） & \kgcsym{同} &  \\
{\sffamily 中3-\allowbreak{}語} & 根号・有理数・無理数・因数・√、∽、全数調査 & \kgcch{2-8}・\kgcch{2-9}・\kgcch{3-10}・\kgcch{4-3}・\kgcch{3-17} & \kgcsym{同} &  \\
{\sffamily 中3-\allowbreak{}扱} & (1)誤差・近似値・$a\times10^{n}$ (2)実数解をもつ二次方程式 (4)円周角の定理の逆 & (1)＝\kgcch{2-6}（誤差・有効数字）、\kgcch{2-7}（10の累乗） & \kgcsym{同}\kgcsep{}\kgcsym{先} & (2)：シリーズは \kgcch{6-8} で虚数解まで進む \\
\end{longtblr}
\Needspace{169pt}\kgbsub{A.10}{高等学校　数学Ⅰ}
\begin{longtblr}[theme=kgc, caption={高等学校　数学Ⅰ}, entry=none, label=none]{colspec={Q[l,wd=19mm] Q[l,wd=31mm] X[2.4,l] Q[l,wd=15mm] X[1.6,l]}, width=\linewidth,
  rowhead=1, row{1}={bg=kgcHead, font=\sffamily\bfseries\footnotesize}, row{even[2]}={bg=kgcBand},
  hlines={0.4pt,bkLine}, vline{2-5}={0.4pt,bkLine}, rowsep=2pt, colsep=3pt, cells={font=\footnotesize}}
番号 & 項目（要点） & 扱っている章 & 扱い & 補足 \\
{\sffamily 数Ⅰ-\allowbreak{}(1)ア(ｱ)} & 実数、無理数の四則 & \kgcch{2-9}、\kgcch{3-10} & \kgcsym{同}\kgcsep{}\kgcsym{先} &  \\
{\sffamily 数Ⅰ-\allowbreak{}(1)ア(ｲ)} & 集合と命題 & \kgcch{4-1}「逆と反例」（仮定・結論） & \kgcsym{一部} & \textbf{集合はなし}（「集合」の語は \kgcch{4-0} の「集合場所」だけ）。「命題」の語も \kgcch{4-1} にない。「必要十分条件」の語は \kgcch{7-5} の専門書の囲みに出る \\
{\sffamily 数Ⅰ-\allowbreak{}(1)ア(ｳ)} & 二次の乗法公式・因数分解 & \kgcch{3-8}、\kgcch{3-10} & \kgcsym{同}\kgcsep{}\kgcsym{先} &  \\
{\sffamily 数Ⅰ-\allowbreak{}(1)ア(ｴ)} & 一次不等式 & \kgcch{3-5}「どこから上になる？：1次不等式の入り口」、\kgcch{2-2}（不等号） & \kgcsym{一部}\kgcsep{}\kgcsym{先} & 「入り口」の扱い \\
{\sffamily 数Ⅰ-\allowbreak{}(2)ア(ｱ)(ｲ)} & 三角比（鋭角・鈍角） & \kgcch{4-6}、\kgcch{4-7}「三角形の名前のつけ方と、鈍角の三角比」 & \kgcsym{同} & 数Ⅰ-扱(3) 0°・90°・180° は \kgcch{4-6}「0°から180°までの三角比」 \\
{\sffamily 数Ⅰ-\allowbreak{}(2)ア(ｳ)} & 正弦定理・余弦定理 & \kgcch{4-7} & \kgcsym{詳} & ロボットの腕の順運動学・逆運動学、力の三角形に使う \\
{\sffamily 数Ⅰ-\allowbreak{}(3)ア(ｱ)(ｲ)} & 二次関数のグラフ・最大最小 & \kgcch{3-9} & \kgcsym{同}\kgcsep{}\kgcsym{先} &  \\
{\sffamily 数Ⅰ-\allowbreak{}(3)ア(ｳ)} & 二次方程式・二次不等式とグラフ & \kgcch{3-10}「式がプラスになるのは、どこ？」 & \kgcsym{同}\kgcsep{}\kgcsym{先} & 本文に「高校では『2次不等式』として学びます」とある \\
{\sffamily 数Ⅰ-\allowbreak{}(4)ア(ｱ)} & 分散・標準偏差・散布図・相関係数 & \kgcch{3-14}、\kgcch{3-15} & \kgcsym{詳}\kgcsep{}\kgcsym{先} & 共分散、$\sigma/\sqrt{n}$（\kgcch{3-14}）、相関係数＝偏差ベクトルの向き（\kgcch{4-12}） \\
{\sffamily 数Ⅰ-\allowbreak{}(4)ア(ｲ)} & コンピュータでの整理 & \kgcch{2-18}（表計算）、\kgcch{3-14}・\kgcch{3-15} & \kgcsym{同} &  \\
{\sffamily 数Ⅰ-\allowbreak{}(4)ア(ｳ)} & 仮説検定の考え方 & \kgcch{3-17} \kgcmotto{}「このコインは、いかさま？（仮説検定）」（\mbox{★★★}） & \kgcsym{一部} & 「もっと知りたい」だけ \\
{\sffamily 数Ⅰ-\allowbreak{}(4)語} & 外れ値 & \kgcch{3-13}、\kgcch{3-15} & \kgcsym{同} &  \\
{\sffamily 数Ⅰ-\allowbreak{}扱} & (2)有限小数・循環小数 & \kgcch{2-9}「分数を小数にすると、終わるかくり返す」 & \kgcsym{同}\kgcsep{}\kgcsym{先} &  \\
\end{longtblr}
\Needspace{199pt}\kgbsub{A.11}{高等学校　数学A}
\begin{longtblr}[theme=kgc, caption={高等学校　数学A}, entry=none, label=none]{colspec={Q[l,wd=19mm] Q[l,wd=31mm] X[2.4,l] Q[l,wd=15mm] X[1.6,l]}, width=\linewidth,
  rowhead=1, row{1}={bg=kgcHead, font=\sffamily\bfseries\footnotesize}, row{even[2]}={bg=kgcBand},
  hlines={0.4pt,bkLine}, vline{2-5}={0.4pt,bkLine}, rowsep=2pt, colsep=3pt, cells={font=\footnotesize}}
番号 & 項目（要点） & 扱っている章 & 扱い & 補足 \\
{\sffamily 数A-\allowbreak{}(1)ア(ｱ)} & 三角形の性質 & \kgcch{4-1}（内角の和）、\kgcch{4-2}（外心）、\kgcch{4-11}（重心）、\kgcch{4-7} \kgcmotto{}「3辺から面積：ヘロンの公式と内接円」 & \kgcsym{一部} & 内心・メネラウス・チェバの語はない \\
{\sffamily 数A-\allowbreak{}(1)ア(ｲ)} & 円の性質 & \kgcch{4-2}（円周角・接線） & \kgcsym{一部} & 方べき・接弦定理の語はない \\
{\sffamily 数A-\allowbreak{}(1)ア(ｳ)} & 空間図形 & \kgcch{4-13} & \kgcsym{一部} & 多面体の語はない \\
{\sffamily 数A-\allowbreak{}(1)イ} & 作図 & なし & \kgcsym{なし} &  \\
{\sffamily 数A-\allowbreak{}(2)ア(ｱ)} & 数え上げの原則 & \kgcch{3-17}（表・樹形図で数える） & \kgcsym{一部} & 和の法則・積の法則の名前はない \\
{\sffamily 数A-\allowbreak{}(2)ア(ｲ)} & 順列・組合せ & \kgcch{3-17}「コインの表の枚数を数える：パスカルの三角形」（組合せ・${}_{n}\mathrm{C}_{r}$）、\kgcch{3-17} \kgcmotto{}「数の爆発：トランプの並べ方」 & \kgcsym{一部} & 「順列」の語・${}_{n}\mathrm{P}_{r}$ はない \\
{\sffamily 数A-\allowbreak{}(2)ア(ｳ)} & 確率・期待値 & \kgcch{3-17}「たくさんやったときの平均：期待値」 & \kgcsym{同} &  \\
{\sffamily 数A-\allowbreak{}(2)ア(ｴ)} & 独立な試行 & なし & \kgcsym{なし} & \kgcch{3-17} \kgcmotto{}で「べつべつにふったサイコロの和では、分散がたし算になる」と言うだけ \\
{\sffamily 数A-\allowbreak{}(2)ア(ｵ)} & 条件付き確率 & なし & \kgcsym{なし} &  \\
{\sffamily 数A-\allowbreak{}(2)語} & ${}_{n}\mathrm{P}_{r}$・${}_{n}\mathrm{C}_{r}$・階乗・n!・排反 & ${}_{n}\mathrm{C}_{r}$＝\kgcch{3-17}（記号カード）、階乗・n!＝\kgcch{5-17} & \kgcsym{一部} & ${}_{n}\mathrm{P}_{r}$・排反はない \\
{\sffamily 数A-\allowbreak{}(3)} & 数学と人間の活動（整数の性質を含む） & \kgcch{2-17}（2進数・16進数）、\kgcch{2-4}（素数・公約数・公倍数）、\kgcch{2-9} \kgcmotto{}（偶数・奇数） & \kgcsym{一部} & 内容の取扱い(4)に約数・倍数、ユークリッドの互除法、二進法、座標の考え方が挙がる（互除法はシリーズにない）。シリーズの \kgcch{2-17} は「2進数」「16進数」の語で扱う \\
\end{longtblr}
\Needspace{169pt}\kgbsub{A.12}{高等学校　数学Ⅱ}
\begin{longtblr}[theme=kgc, caption={高等学校　数学Ⅱ}, entry=none, label=none]{colspec={Q[l,wd=19mm] Q[l,wd=31mm] X[2.4,l] Q[l,wd=15mm] X[1.6,l]}, width=\linewidth,
  rowhead=1, row{1}={bg=kgcHead, font=\sffamily\bfseries\footnotesize}, row{even[2]}={bg=kgcBand},
  hlines={0.4pt,bkLine}, vline{2-5}={0.4pt,bkLine}, rowsep=2pt, colsep=3pt, cells={font=\footnotesize}}
番号 & 項目（要点） & 扱っている章 & 扱い & 補足 \\
{\sffamily 数Ⅱ-\allowbreak{}(1)ア(ｱ)} & 三次の乗法公式・因数分解 & \kgcch{3-8}「立方体で見る (a＋b)$^{3}$」 & \kgcsym{一部}\kgcsep{}\kgcsym{先} & 三次の因数分解の公式はない \\
{\sffamily 数Ⅱ-\allowbreak{}(1)ア(ｲ)} & 多項式の除法・分数式 & \kgcch{9-7}（部分分数の前に「2s＋5＝2(s＋2)＋1」の1行） & \kgcsym{一部} & 除法の節はない \\
{\sffamily 数Ⅱ-\allowbreak{}(1)ア(ｳ)} & 複素数の四則 & \kgcch{6-8} & \kgcsym{同}\kgcsep{}\kgcsym{後} & 虚数単位は \textbf{j}（i は電流に使う） \\
{\sffamily 数Ⅱ-\allowbreak{}(1)ア(ｴ)} & 判別、解と係数の関係 & \kgcch{3-10}（判別式）、\kgcch{6-8}「2次方程式の、虚数の解」、\kgcch{6-12}（解と係数の関係） & \kgcsym{同} &  \\
{\sffamily 数Ⅱ-\allowbreak{}(1)ア(ｵ)} & 因数定理・高次方程式 & \kgcch{9-9}「3次の特性方程式：強くしすぎると、こわれる」、\kgcch{6-8} \kgcmotto{}「複素数まで広げれば、…」 & \kgcsym{一部} & 因数定理・剰余の定理の語はない \\
{\sffamily 数Ⅱ-\allowbreak{}(1)イ} & 等式・不等式の証明 & なし & \kgcsym{なし} & 恒等式・係数比較は \kgcch{3-8} にある \\
{\sffamily 数Ⅱ-\allowbreak{}(1)語} & 二項定理・虚数・i & 二項定理＝\kgcch{3-8} \kgcmotto{}「係数のならび：1・2・1、1・3・3・1」（パスカルの三角形）、\kgcch{5-4} \kgcmotto{}「どんな n でも成り立つわけ」（(x＋h)$^{n}$ の係数）、虚数＝\kgcch{6-8} & \kgcsym{一部} & 「二項定理」の語はない。i ではなく j \\
{\sffamily 数Ⅱ-\allowbreak{}(2)ア(ｱ)} & 内分点・外分点・2点間の距離 & \kgcch{4-11}（内分点）、\kgcch{4-4}「座標の上の2点の距離」 & \kgcsym{一部} & 外分点はない \\
{\sffamily 数Ⅱ-\allowbreak{}(2)ア(ｲ)} & 直線と円の方程式 & \kgcch{3-5}（2点・1点と傾きから直線の式）、\kgcch{4-4}（$(x-a)^{2}+(y-b)^{2}$\allowbreak$=r^{2}$）、\kgcch{4-12} \kgcmotto{}「点と直線の距離、直交する2直線」 & \kgcsym{同}\kgcsep{}\kgcsym{先} &  \\
{\sffamily 数Ⅱ-\allowbreak{}(2)ア(ｳ)} & 軌跡 & なし & \kgcsym{なし} & \kgcch{9-9} の「根軌跡」は制御の用語で別もの \\
{\sffamily 数Ⅱ-\allowbreak{}(2)ア(ｴ)} & 不等式の表す領域 & なし & \kgcsym{なし} &  \\
{\sffamily 数Ⅱ-\allowbreak{}(3)ア(ｱ)} & 指数の拡張（整数→\nobreak{}有理数） & \kgcch{2-7}「0乗とマイナスの指数」、\kgcch{2-12}、\kgcch{3-11}「分数の指数」 & \kgcsym{同}\kgcsep{}\kgcsym{先} & 0乗・負の指数は第2巻（\mbox{★〜★★}） \\
{\sffamily 数Ⅱ-\allowbreak{}(3)ア(ｲ)〜(ｴ)} & 指数関数・対数・対数関数 & \kgcch{3-11}、\kgcch{3-12} & \kgcsym{詳}\kgcsep{}\kgcsym{先} & 片対数グラフ・デシベル・底の変換（電卓）まで \\
{\sffamily 数Ⅱ-\allowbreak{}(3)語} & 累乗根・$\log_{a}x$・常用対数 & \kgcch{3-11}、\kgcch{3-12} & \kgcsym{同} &  \\
{\sffamily 数Ⅱ-\allowbreak{}(4)ア(ｱ)} & 一般角・弧度法 & \kgcch{4-8}「何周しても使える角：一般角」、\kgcch{4-2}「角を長さではかる：弧度法とラジアン」 & \kgcsym{同} &  \\
{\sffamily 数Ⅱ-\allowbreak{}(4)ア(ｲ)(ｳ)} & 三角関数のグラフ・相互関係 & \kgcch{4-8} & \kgcsym{同} &  \\
{\sffamily 数Ⅱ-\allowbreak{}(4)ア(ｴ)} & 加法定理・2倍角・合成 & \kgcch{4-9} & \kgcsym{詳} & 加法定理を2通りで導く、積と和の式、うなり、sin$^{2}$ の平均 \\
{\sffamily 数Ⅱ-\allowbreak{}(5)ア(ｱ)} & 微分係数・導関数 & \kgcch{5-1}〜\kgcch{5-4} & \kgcsym{詳} & 数Ⅱ-扱(1)は「三次までの関数を中心」。シリーズは \kgcch{5-4} で x$^{n}$ 一般 \\
{\sffamily 数Ⅱ-\allowbreak{}(5)ア(ｲ)} & 増減・極大極小 & \kgcch{5-7} & \kgcsym{同} &  \\
{\sffamily 数Ⅱ-\allowbreak{}(5)ア(ｳ)} & 不定積分・定積分 & \kgcch{5-9}〜\kgcch{5-11} & \kgcsym{詳} & 数Ⅱ-扱(1)は「二次までの関数を中心」。区分求積（\kgcch{5-10}）から入る \\
{\sffamily 数Ⅱ-\allowbreak{}(5)イ} & 面積 & \kgcch{5-9}、\kgcch{5-11}「横軸より下の部分」 & \kgcsym{同} &  \\
{\sffamily 数Ⅱ-\allowbreak{}(5)語} & 極限値・lim & \kgcch{5-3}（記号カード lim） & \kgcsym{同} & 極限は数値で近づく様子で（$\varepsilon$-$\delta$は使わない） \\
\end{longtblr}
\Needspace{199pt}\kgbsub{A.13}{高等学校　数学B}
\begin{longtblr}[theme=kgc, caption={高等学校　数学B}, entry=none, label=none]{colspec={Q[l,wd=19mm] Q[l,wd=31mm] X[2.4,l] Q[l,wd=15mm] X[1.6,l]}, width=\linewidth,
  rowhead=1, row{1}={bg=kgcHead, font=\sffamily\bfseries\footnotesize}, row{even[2]}={bg=kgcBand},
  hlines={0.4pt,bkLine}, vline{2-5}={0.4pt,bkLine}, rowsep=2pt, colsep=3pt, cells={font=\footnotesize}}
番号 & 項目（要点） & 扱っている章 & 扱い & 補足 \\
{\sffamily 数B-\allowbreak{}(1)ア(ｱ)} & 等差数列・等比数列 & \kgcch{7-10}（等比数列・等比級数） & \kgcsym{一部} & \textbf{等差数列はない}（\kgcch{6-7} で「等差」の語が1回出るだけ） \\
{\sffamily 数B-\allowbreak{}(1)ア(ｲ)} & いろいろな数列の和 & \kgcch{5-10}「たし算をまとめて書く：$\Sigma$」（$1^{2}+2^{2}+\dots+n^{2}$\allowbreak$=n(n+1)(2n+1)/6$） & \kgcsym{一部} &  \\
{\sffamily 数B-\allowbreak{}(1)ア(ｳ)} & 漸化式 & \kgcch{6-2}（オイラー法 $y_{n+1}$＝…の形。「漸化式」の語はない）、\kgcch{7-10}「行列のくり返し：漸化式とオイラー法を行列で」 & \kgcsym{一部} &  \\
{\sffamily 数B-\allowbreak{}(1)ア(ｴ)} & 数学的帰納法 & なし & \kgcsym{なし} & \kgcch{5-4} \kgcmotto{}「どんな n でも成り立つわけ」は (x＋h)$^{n}$ の係数（パスカルの三角形）による説明で、帰納法ではない \\
{\sffamily 数B-\allowbreak{}(1)語} & $\Sigma$ & \kgcch{5-10} & \kgcsym{同} &  \\
{\sffamily 数B-\allowbreak{}(2)ア(ｱ)} & 標本調査の考え方 & \kgcch{3-17}、\kgcch{3-14}「何回もはかって平均すると、ずれは $\sigma/\sqrt{n}$」 & \kgcsym{一部} &  \\
{\sffamily 数B-\allowbreak{}(2)ア(ｲ)} & 確率変数と確率分布 & \kgcch{3-17}（コインの表の枚数の分布、\kgcmotto{}でサイコロの和の分散）、\kgcch{6-4}「めったに起きないことの待ち時間」（指数分布・ポアソン分布） & \kgcsym{一部} & 「確率変数」「確率分布」の語はない \\
{\sffamily 数B-\allowbreak{}(2)ア(ｳ)} & 二項分布・正規分布 & \kgcch{3-17}「たくさん足すと、つりがね形：正規分布」（±$\sigma$ に約68 \%、±2$\sigma$ に約95 \%）、\kgcch{3-17} のパスカルの三角形 & \kgcsym{一部} & 「二項分布」の語はない \\
{\sffamily 数B-\allowbreak{}(2)ア(ｴ)} & 区間推定・仮説検定 & 仮説検定＝\kgcch{3-17} \kgcmotto{}（\mbox{★★★}）。区間推定＝なし & \kgcsym{一部}\kgcor{}\kgcsym{なし} &  \\
{\sffamily 数B-\allowbreak{}(2)語} & 信頼区間・有意水準 & なし & \kgcsym{なし} &  \\
{\sffamily 数B-\allowbreak{}(3)} & 数学と社会生活（扱(3)散布図のデータを関数とみなす） & \kgcch{3-16}（最小二乗法・回帰直線） & \kgcsym{詳} &  \\
\end{longtblr}
\Needspace{229pt}\kgbsub{A.14}{高等学校　数学Ⅲ}
\begin{longtblr}[theme=kgc, caption={高等学校　数学Ⅲ}, entry=none, label=none]{colspec={Q[l,wd=19mm] Q[l,wd=31mm] X[2.4,l] Q[l,wd=15mm] X[1.6,l]}, width=\linewidth,
  rowhead=1, row{1}={bg=kgcHead, font=\sffamily\bfseries\footnotesize}, row{even[2]}={bg=kgcBand},
  hlines={0.4pt,bkLine}, vline{2-5}={0.4pt,bkLine}, rowsep=2pt, colsep=3pt, cells={font=\footnotesize}}
番号 & 項目（要点） & 扱っている章 & 扱い & 補足 \\
{\sffamily 数Ⅲ-\allowbreak{}(1)ア(ｱ)} & 数列の極限・$\{r^{n}\}$ & \kgcch{5-10}「n を大きくした先を、式で求める」、\kgcch{5-14}「$(1+1/n)^{n}$」、\kgcch{7-10}（$\lambda^{n}$\allowbreak$=r^{n}e^{jn\theta}$、固有値の大きさで「ふえる・消える」） & \kgcsym{一部} & 数列の極限の節はない \\
{\sffamily 数Ⅲ-\allowbreak{}(1)ア(ｲ)} & 無限級数・無限等比級数 & \kgcch{7-10}「くり返しの合計：等比級数」（$1+x+x^{2}+\dots$\allowbreak$=1/(1-x)$）、\kgcch{5-17}（テイラー展開）、\kgcch{9-11}「無限に続くたし算で書く」 & \kgcsym{一部} &  \\
{\sffamily 数Ⅲ-\allowbreak{}(1)ア(ｳ)} & 分数関数・無理関数 & \kgcch{3-7}（y＝a/x）、\kgcch{3-7} \kgcmotto{}（1/x$^{2}$） & \kgcsym{一部} & 平行移動した分数関数、無理関数のグラフはない \\
{\sffamily 数Ⅲ-\allowbreak{}(1)ア(ｴ)} & 合成関数・逆関数 & \kgcch{3-3}「箱をつなぐ：合成関数」、\kgcch{3-12}（逆関数：本文の囲み）、\kgcch{5-14}（ln） & \kgcsym{同}\kgcsep{}\kgcsym{先} &  \\
{\sffamily 数Ⅲ-\allowbreak{}(1)ア(ｵ)} & 関数の値の極限 & \kgcch{5-3}「かぎりなく近づく：極限と lim」 & \kgcsym{同} &  \\
{\sffamily 数Ⅲ-\allowbreak{}(2)ア(ｱ)} & 微分可能性・積・商 & \kgcch{5-6}（積の微分）、\kgcch{5-6} \kgcmotto{}「商の微分」、\kgcch{5-3} \kgcmotto{}「とがった所には、接線が引けない」 & \kgcsym{同}\newline{}（商は一部） &  \\
{\sffamily 数Ⅲ-\allowbreak{}(2)ア(ｲ)} & 合成関数の導関数 & \kgcch{5-6} & \kgcsym{同} &  \\
{\sffamily 数Ⅲ-\allowbreak{}(2)ア(ｳ)} & 三角・指数・対数関数の導関数 & \kgcch{5-14}、\kgcch{5-15} & \kgcsym{同} &  \\
{\sffamily 数Ⅲ-\allowbreak{}(2)ア(ｴ)} & 接線・増減・凹凸・グラフ & \kgcch{5-3}「接線の式」、\kgcch{5-7}「2回微分で、山と谷を見分ける」 & \kgcsym{一部} & 「変曲点」の語はない。「下に凸」は \kgcch{5-17} に出る \\
{\sffamily 数Ⅲ-\allowbreak{}(2)語} & 自然対数・e・変曲点 & \kgcch{5-14}（e・ln）、\kgcch{5-5}「2回微分する」 & \kgcsym{一部} & 変曲点なし。第二次導関数は〔用語・記号〕ではなく(2)イ(ｲ)の本文にある \\
{\sffamily 数Ⅲ-\allowbreak{}(3)ア(ｱ)} & 不定積分・定積分の性質 & \kgcch{5-11} & \kgcsym{同} &  \\
{\sffamily 数Ⅲ-\allowbreak{}(3)ア(ｲ)} & 置換積分・部分積分 & \kgcch{5-16} & \kgcsym{同} &  \\
{\sffamily 数Ⅲ-\allowbreak{}(3)ア(ｳ)} & 面積・体積・曲線の長さ & 面積＝\kgcch{5-11}。体積＝\kgcch{5-10} の例題「うすい板に分けて体積を求める」（円すい）と章末問題（四角すい）。曲線の長さ＝\kgcch{6-11} のせいかくに言うと（$s=\int\!\sqrt{1+y'^{2}}\,dx$）、\kgcch{8-7}「長さでたす線積分」 & \kgcsym{一部} & 回転体の体積の計算演習・曲線の長さの節はない \\
{\sffamily 数Ⅲ-\allowbreak{}扱(1)} & 直線上・平面上の点の速度・加速度 & \kgcch{5-5}、\kgcch{5-15}「円の上を回る点の速さと向き」、\kgcch{4-11} & \kgcsym{詳} &  \\
{\sffamily 数Ⅲ-\allowbreak{}注} & 微分方程式（本文になし） & 第6巻全体 & \kgcsym{範囲外} &  \\
\end{longtblr}
\Needspace{169pt}\kgbsub{A.15}{高等学校　数学C}
\begin{longtblr}[theme=kgc, caption={高等学校　数学C}, entry=none, label=none]{colspec={Q[l,wd=19mm] Q[l,wd=31mm] X[2.4,l] Q[l,wd=15mm] X[1.6,l]}, width=\linewidth,
  rowhead=1, row{1}={bg=kgcHead, font=\sffamily\bfseries\footnotesize}, row{even[2]}={bg=kgcBand},
  hlines={0.4pt,bkLine}, vline{2-5}={0.4pt,bkLine}, rowsep=2pt, colsep=3pt, cells={font=\footnotesize}}
番号 & 項目（要点） & 扱っている章 & 扱い & 補足 \\
{\sffamily 数C-\allowbreak{}(1)ア(ｱ)} & ベクトル・位置ベクトル・成分 & \kgcch{4-10}、\kgcch{4-11} & \kgcsym{同} & \kgcch{4-10} は\mbox{★★}、\kgcch{4-11} は\mbox{★★〜★★★} \\
{\sffamily 数C-\allowbreak{}(1)ア(ｲ)} & 内積 & \kgcch{4-12} & \kgcsym{詳} & 仕事・正射影・相関係数 \\
{\sffamily 数C-\allowbreak{}(1)ア(ｳ)} & 空間への拡張 & \kgcch{4-13} & \kgcsym{詳} & 法線ベクトル・空間の直線と平面・右手系。外積は \kgcch{4-14}（範囲外） \\
{\sffamily 数C-\allowbreak{}(2)ア(ｱ)} & 放物線・楕円・双曲線（二次曲線） & 放物線＝\kgcch{3-9}（パラボラアンテナの焦点）。だ円＝\kgcch{4-8} \kgcmotto{}「円をつぶすと、だ円」（\mbox{★★★★}）、\kgcch{4-15}（極方程式）、\kgcch{8-1}「等高線の式は、だ円」、\kgcch{7-9} \kgcmotto{}。双曲線＝\kgcch{3-7}（y＝a/x のグラフとして） & \kgcsym{一部} & 焦点・準線で3つをまとめる扱いはない。語は「楕円」でなく「だ円」 \\
{\sffamily 数C-\allowbreak{}(2)ア(ｲ)} & 媒介変数表示 & \kgcch{4-8}（だ円の媒介変数）、\kgcch{4-11}（$\vec{p}=\vec{p}_{0}+t\vec{v}$）、\kgcch{8-7}「道を1つの数 t で表す」 & \kgcsym{同} &  \\
{\sffamily 数C-\allowbreak{}(2)ア(ｳ)} & 極座標・極方程式 & \kgcch{4-2}「距離と向きで場所を言う：極座標」、\kgcch{4-15}「曲線を、極座標の式で表す」 & \kgcsym{同} &  \\
{\sffamily 数C-\allowbreak{}(2)ア(ｴ)} & 複素数平面・極形式・四則の図形的意味 & \kgcch{6-9}「複素数を点で表す：複素平面」「向きの角：偏角と極形式」 & \kgcsym{同} & 語は「複素平面」 \\
{\sffamily 数C-\allowbreak{}(2)ア(ｵ)} & ド・モアブルの定理 & \kgcch{6-9}「何回もかける：ド・モアブルの定理」 & \kgcsym{同} & \kgcmotto{}「1の n 乗根」 \\
{\sffamily 数C-\allowbreak{}(3)ア(ｱ)} & 図・表・統計グラフ & \kgcch{1-17}、\kgcch{1-18}、\kgcch{3-4}、\kgcch{3-13} & \kgcsym{同} &  \\
{\sffamily 数C-\allowbreak{}(3)ア(ｲ)} & 離散グラフや行列で表現する意義 & 行列＝第7巻全体（\kgcch{3-6}「係数を表にならべる：行列の種」、\kgcch{1-13} が種） & \kgcsym{詳}\newline{}（行列） & 行列は演算・行列式・逆行列・固有値・対角化まで（指導要領の範囲外） \\
\end{longtblr}
\bkevenpage\kgbsec{B}{付録B　全152章の一覧}\label{kgc-b}\label{kgc-sec11}
この付録は、全9巻・152章（序章9・章139・読みもの3・終章1）を、巻と章の順に並べた一覧です。章名・★・路線は各巻のもくじと同じです。「指導要領の学年・科目」は、その章の内容が当たる学習指導要領の学年や科目で、1つの章が複数の学年にまたがるときは「入口〜出口」（最も低い学年〜最も高い学年・科目）で示しています。

\begin{kgclegend}{表の読み方}
\begin{tblr}{colspec={Q[l,wd=24mm] X[l]}, rowsep=2pt, colsep=3pt, cells={font=\small, valign=m}}
★ & 章の内容の高さの目安です。塗りの★はその章の下の高さ、白抜きの☆は上の高さまでを表します（例：\kgcstars{2}{3}＝★★〜★★★）。★は指導要領の学年と一対一ではありません。 \\
路線 & 7本の路線のどの駅に当たるかを、形の記号と番号で示しています。\par\smallskip{\raggedright \mbox{\kgcroute{1}\,変化（ひき算→傾き→微分）}\hspace{1.2em plus 1fill}\mbox{\kgcroute{2}\,たし合わせ（かけ算→面積→積分）}\hspace{1.2em plus 1fill}\mbox{\kgcroute{3}\,倍々（累乗→指数→e）}\hspace{1.2em plus 1fill}\mbox{\kgcroute{4}\,回転と波（角→sin→複素数）}\hspace{1.2em plus 1fill}\mbox{\kgcroute{5}\,表（九九の表→行列）}\hspace{1.2em plus 1fill}\mbox{\kgcroute{6}\,矢印（向き→ベクトル→場）}\hspace{1.2em plus 1fill}\mbox{\kgcroute{7}\,はかる（目もり→誤差→統計）}\par} \\
指導要領の学年・科目 & 小＝小学校、中＝中学校、数Ⅰ〜数C＝高等学校の科目（数学Ⅰ〜数学C）、大学＝指導要領の範囲外で、大学の工学部で一般に学ぶ内容です。「範囲外」は、指導要領にも大学の科目の一覧にも項目がない内容です。「―」は、特定の学年に当てはめていない章です（1-0・2-0・2-18・2-19・4-0・5-0・9-13）。 \\
範囲外 & \kgcout{a}＝指導要領の範囲外の内容が中心の章、\kgcout{p}＝一部に指導要領の範囲外の内容を含む章、印なし＝指導要領の範囲の中の章です。 \\
\end{tblr}
\end{kgclegend}

\Needspace{8\baselineskip}\kgbsub{第1巻}{かぞえる・はかる}\par\nobreak{\small\noindent 学年の目安：小学1年〜3年　章数：22}\par\nobreak\vspace{3pt}
\begin{longtblr}[theme=kgc, caption={第1巻　かぞえる・はかる}, entry=none, label=none]{colspec={Q[l,wd=14mm] X[l] Q[l,wd=17mm] Q[l,wd=12.5mm] Q[l,wd=25mm] Q[c,wd=10mm]}, width=\linewidth,
  rowhead=1, row{1}={bg=kgcHead, font=\sffamily\bfseries\footnotesize}, row{even[2]}={bg=kgcBand},
  hlines={0.4pt,bkLine}, vline{2-6}={0.4pt,bkLine}, rowsep=1.5pt, colsep=3pt, cells={font=\footnotesize, valign=m}, row{1}={valign=m}}
章 & 章名 & ★ & 路線 & 指導要領の学年・科目 & 範囲外 \\
{\sffamily\bfseries 1-0}\newline{\footnotesize\sffamily 序章} & かずを つかうと、なにが いいの？ & \kgcstars{1}{1} & \kgcroute{7} & ― &  \\
{\sffamily\bfseries 1-1} & いくつ あるかな？ & \kgcstars{1}{1} & \kgcroute{2} & 小1 &  \\
{\sffamily\bfseries 1-2} & 10の まとまりを つくろう & \kgcstars{1}{1} & \kgcroute{3} & 小1〜小3 &  \\
{\sffamily\bfseries 1-3} & なにも ないときの「0」 & \kgcstars{1}{1} & \kgcroute{6} & 小1〜小2 &  \\
{\sffamily\bfseries 1-4} & たし算は「あわせる」「ふえる」 & \kgcstars{1}{1} & \kgcroute{2} & 小1 &  \\
{\sffamily\bfseries 1-5} & ひき算は「のこり」「ちがい」「かわった 量」 & \kgcstars{1}{1} & \kgcroute{1} & 小1 &  \\
{\sffamily\bfseries 1-6} & かずの 線：数直線 & \kgcstars{1}{1} & \kgcroute{6} & 小1〜小3 &  \\
{\sffamily\bfseries 1-7} & 大きな かずの たし算・ひき算 & \kgcstars{1}{1} & \kgcroute{2} & 小1〜小3 &  \\
{\sffamily\bfseries 1-8} & ながさを はかろう & \kgcstars{1}{1} & \kgcroute{7} & 小1〜小2 &  \\
{\sffamily\bfseries 1-9} & かさと おもさ & \kgcstars{1}{1} & \kgcroute{7} & 小1〜小3 &  \\
{\sffamily\bfseries 1-10} & 時こくと 時間 & \kgcstars{1}{1} & \kgcroute{1} & 小1〜小3 &  \\
{\sffamily\bfseries 1-11} & 目もりを よもう & \kgcstars{1}{1} & \kgcroute{7} & 小2〜小3 &  \\
{\sffamily\bfseries 1-12} & かけ算は「同じ かずの まとまり」 & \kgcstars{1}{1} & \kgcroute{2} & 小2 &  \\
{\sffamily\bfseries 1-13} & 九九の 表を たんけんしよう & \kgcstars{1}{1} & \kgcroute{5} & 小2〜小3 &  \\
{\sffamily\bfseries 1-14} & わり算は「同じ かずずつ 分ける」 & \kgcstars{1}{1} & \kgcroute{1} & 小3 &  \\
{\sffamily\bfseries 1-15} & あまりの ある わり算 & \kgcstars{1}{1} & \kgcroute{4} & 小3〜小4 &  \\
{\sffamily\bfseries 1-16} & 1より 小さい かず：小数と 分数 & \kgcstars{1}{1} & \kgcroute{7} & 小2〜小5 &  \\
{\sffamily\bfseries 1-17} & 表と ぼうグラフ & \kgcstars{1}{1} & \kgcroute{5} & 小1〜小3 &  \\
{\sffamily\bfseries 1-18} & おれ線グラフ：かわり方を 見る & \kgcstars{1}{1} & \kgcroute{1} & 小4 &  \\
{\sffamily\bfseries 1-19} & かたち：三角・四角・円 & \kgcstars{1}{1} & \kgcroute{4} & 小1〜小3 &  \\
{\sffamily\bfseries 1-20} & 角の 大きさ & \kgcstars{1}{1} & \kgcroute{4} & 小3〜小4 &  \\
{\sffamily\bfseries 1-21} & 広さを くらべよう：面積 & \kgcstars{1}{1} & \kgcroute{2} & 小1〜小4 &  \\
\end{longtblr}
\Needspace{8\baselineskip}\kgbsub{第2巻}{数と量のことば}\par\nobreak{\small\noindent 学年の目安：小学4年〜中学1年　章数：20}\par\nobreak\vspace{3pt}
\begin{longtblr}[theme=kgc, caption={第2巻　数と量のことば}, entry=none, label=none]{colspec={Q[l,wd=14mm] X[l] Q[l,wd=17mm] Q[l,wd=12.5mm] Q[l,wd=25mm] Q[c,wd=10mm]}, width=\linewidth,
  rowhead=1, row{1}={bg=kgcHead, font=\sffamily\bfseries\footnotesize}, row{even[2]}={bg=kgcBand},
  hlines={0.4pt,bkLine}, vline{2-6}={0.4pt,bkLine}, rowsep=1.5pt, colsep=3pt, cells={font=\footnotesize, valign=m}, row{1}={valign=m}}
章 & 章名 & ★ & 路線 & 指導要領の学年・科目 & 範囲外 \\
{\sffamily\bfseries 2-0}\newline{\footnotesize\sffamily 序章} & 「だんだん速く」を数で言うと？ & \kgcstars{1}{1} & \kgcroute{1} & ― &  \\
{\sffamily\bfseries 2-1} & 「かぞえる数」と「はかる数」 & \kgcstars{1}{1} & \kgcroute{7} & 小3〜小5 &  \\
{\sffamily\bfseries 2-2} & マイナスの数と数直線 & \kgcstars{1}{2} & \kgcroute{6} & 小2〜中1 &  \\
{\sffamily\bfseries 2-3} & ＋ － × ÷ は、何を表しているの？ & \kgcstars{1}{1} & \kgcroute{1}\par \kgcroute{2} & 小1〜小5 &  \\
{\sffamily\bfseries 2-4} & 分数と小数の計算 & \kgcstars{1}{2} & \kgcroute{1} & 小4〜中1 &  \\
{\sffamily\bfseries 2-5} & 比・割合・百分率 & \kgcstars{1}{2} & \kgcroute{7} & 小4〜小6 &  \\
{\sffamily\bfseries 2-6} & はかった数には「ずれ」がある & \kgcstars{1}{2} & \kgcroute{7} & 小4〜中3 &  \\
{\sffamily\bfseries 2-7} & 何回もかける：累乗と指数 & \kgcstars{1}{2} & \kgcroute{3} & 小4〜数Ⅱ &  \\
{\sffamily\bfseries 2-8} & 2回かけると□になる数：平方根 & \kgcstars{2}{2} & \kgcroute{3} & 中3 &  \\
{\sffamily\bfseries 2-9} & 分数で書けない数がある？ & \kgcstars{2}{2} & \kgcroute{6} & 小5〜数Ⅰ &  \\
{\sffamily\bfseries 2-10} & 単位の世界ルール：SI単位 & \kgcstars{1}{2} & \kgcroute{7} & 小2〜小5 & \kgcout{p} \\
{\sffamily\bfseries 2-11} & なぜ文字を使うの？ & \kgcstars{2}{2} & \kgcroute{5} & 小2〜中1 &  \\
{\sffamily\bfseries 2-12} & 文字式のやくそく & \kgcstars{2}{2} & \kgcroute{5} & 中1〜数Ⅱ &  \\
{\sffamily\bfseries 2-13} & 公式に数を入れてみよう：代入 & \kgcstars{2}{2} & \kgcroute{5} & 小4〜中1 &  \\
{\sffamily\bfseries 2-14} & 式を変形して、ほしい量をとり出そう & \kgcstars{2}{2} & \kgcroute{5} & 中1〜中2 &  \\
{\sffamily\bfseries 2-15} & 比例と反比例 & \kgcstars{1}{2} & \kgcroute{1} & 小5〜中1 &  \\
{\sffamily\bfseries 2-16} & 何秒後に追いつく？1次方程式 & \kgcstars{2}{2} & \kgcroute{1} & 中1 &  \\
{\sffamily\bfseries 2-17} & コンピュータの数：2進数と16進数 & \kgcstars{2}{2} & \kgcroute{3} & 数A & \kgcout{p} \\
{\sffamily\bfseries 2-18} & 表計算で計算しよう & \kgcstars{2}{2} & \kgcroute{5} & ― & \kgcout{p} \\
{\sffamily\bfseries 2-19} & プログラムで計算しよう & \kgcstars{2}{2} & \kgcroute{1} & ― & \kgcout{p} \\
\end{longtblr}
\Needspace{8\baselineskip}\kgbsub{第3巻}{関数とデータ　きまりを見つける}\par\nobreak{\small\noindent 学年の目安：中学1年〜3年　章数：18}\par\nobreak\vspace{3pt}
\begin{longtblr}[theme=kgc, caption={第3巻　関数とデータ}, entry=none, label=none]{colspec={Q[l,wd=14mm] X[l] Q[l,wd=17mm] Q[l,wd=12.5mm] Q[l,wd=25mm] Q[c,wd=10mm]}, width=\linewidth,
  rowhead=1, row{1}={bg=kgcHead, font=\sffamily\bfseries\footnotesize}, row{even[2]}={bg=kgcBand},
  hlines={0.4pt,bkLine}, vline{2-6}={0.4pt,bkLine}, rowsep=1.5pt, colsep=3pt, cells={font=\footnotesize, valign=m}, row{1}={valign=m}}
章 & 章名 & ★ & 路線 & 指導要領の学年・科目 & 範囲外 \\
{\sffamily\bfseries 3-0}\newline{\footnotesize\sffamily 序章} & はかったデータから、きまりを見つけよう & \kgcstars{2}{2} & \kgcroute{7} & 中1 &  \\
{\sffamily\bfseries 3-1} & 変わる量と変わらない量 & \kgcstars{1}{1} & \kgcroute{7} & 中1 &  \\
{\sffamily\bfseries 3-2} & 場所を数で表す：座標 & \kgcstars{1}{2} & \kgcroute{6} & 小4〜中1 &  \\
{\sffamily\bfseries 3-3} & 関数は「入れると出てくる」ルール & \kgcstars{1}{2} & \kgcroute{5} & 中1〜数Ⅲ & \kgcout{p} \\
{\sffamily\bfseries 3-4} & グラフをかこう・読もう & \kgcstars{1}{1} & \kgcroute{7} & 中1 &  \\
{\sffamily\bfseries 3-5} & グラフの傾きは「変化の速さ」 & \kgcstars{2}{2} & \kgcroute{1} & 中2〜数Ⅱ &  \\
{\sffamily\bfseries 3-6} & 2本の直線が交わるところ：連立方程式 & \kgcstars{2}{2} & \kgcroute{5} & 中2 &  \\
{\sffamily\bfseries 3-7} & 反比例のグラフ & \kgcstars{2}{2} & \kgcroute{1} & 小6〜数Ⅲ &  \\
{\sffamily\bfseries 3-8} & かっこをはずす：展開と平方完成 & \kgcstars{2}{3} & \kgcroute{1} & 中3〜数Ⅱ &  \\
{\sffamily\bfseries 3-9} & 放物線と $y=ax^2$ & \kgcstars{2}{3} & \kgcroute{1} & 中3〜数Ⅰ &  \\
{\sffamily\bfseries 3-10} & ボールはいつその高さに届く？：2次方程式 & \kgcstars{2}{3} & \kgcroute{1} & 中3〜数Ⅱ &  \\
{\sffamily\bfseries 3-11} & 指数関数：いまある量に比例してふえる・へる & \kgcstars{2}{3} & \kgcroute{3} & 数Ⅱ &  \\
{\sffamily\bfseries 3-12} & 対数：けたを数える道具 & \kgcstars{2}{3} & \kgcroute{3} & 数Ⅱ〜数Ⅲ & \kgcout{p} \\
{\sffamily\bfseries 3-13} & データのちらばり：ヒストグラムと平均 & \kgcstars{2}{2} & \kgcroute{7} & 小6〜数Ⅰ &  \\
{\sffamily\bfseries 3-14} & ばらつきを1つの数で：標準偏差 & \kgcstars{3}{3} & \kgcroute{7} & 数Ⅰ〜数B &  \\
{\sffamily\bfseries 3-15} & 2つの量の関係：散布図と相関 & \kgcstars{3}{3} & \kgcroute{7} & 数Ⅰ &  \\
{\sffamily\bfseries 3-16} & データに直線を引く：最小二乗法 & \kgcstars{3}{3} & \kgcroute{7} & 数B〜大学 & \kgcout{p} \\
{\sffamily\bfseries 3-17} & 偶然のばらつき：確率と正規分布 & \kgcstars{3}{3} & \kgcroute{7} & 小6〜数B &  \\
\end{longtblr}
\Needspace{8\baselineskip}\kgbsub{第4巻}{形と向き　図形・三角関数・ベクトル}\par\nobreak{\small\noindent 学年の目安：中学2年〜高校2年　章数：16}\par\nobreak\vspace{3pt}
\begin{longtblr}[theme=kgc, caption={第4巻　形と向き}, entry=none, label=none]{colspec={Q[l,wd=14mm] X[l] Q[l,wd=17mm] Q[l,wd=12.5mm] Q[l,wd=25mm] Q[c,wd=10mm]}, width=\linewidth,
  rowhead=1, row{1}={bg=kgcHead, font=\sffamily\bfseries\footnotesize}, row{even[2]}={bg=kgcBand},
  hlines={0.4pt,bkLine}, vline{2-6}={0.4pt,bkLine}, rowsep=1.5pt, colsep=3pt, cells={font=\footnotesize, valign=m}, row{1}={valign=m}}
章 & 章名 & ★ & 路線 & 指導要領の学年・科目 & 範囲外 \\
{\sffamily\bfseries 4-0}\newline{\footnotesize\sffamily 序章} & ロボットの手は、どこにある？ & \kgcstars{2}{2} & \kgcroute{4} & ― & \kgcout{p} \\
{\sffamily\bfseries 4-1} & きまりから考えを積み上げる：定義・公理・定理 & \kgcstars{2}{3} & \kgcroute{6} & 小3〜数Ⅰ &  \\
{\sffamily\bfseries 4-2} & 角度と円、ラジアン & \kgcstars{1}{2} & \kgcroute{4} & 小3〜数C &  \\
{\sffamily\bfseries 4-3} & にた形と比：相似 & \kgcstars{2}{3} & \kgcroute{6} & 小6〜中3 &  \\
{\sffamily\bfseries 4-4} & 三平方の定理 & \kgcstars{3}{3} & \kgcroute{6} & 中3〜数Ⅱ &  \\
{\sffamily\bfseries 4-5} & 面積と体積：細かく分けてたす & \kgcstars{1}{2} & \kgcroute{2} & 小4〜中3 &  \\
{\sffamily\bfseries 4-6} & 坂道と三角比 & \kgcstars{2}{3} & \kgcroute{4} & 数Ⅰ & \kgcout{p} \\
{\sffamily\bfseries 4-7} & 三角形を解く：正弦定理と余弦定理 & \kgcstars{3}{3} & \kgcroute{4} & 数Ⅰ〜数A & \kgcout{p} \\
{\sffamily\bfseries 4-8} & 三角関数：円運動の影 & \kgcstars{3}{4} & \kgcroute{4} & 数Ⅱ〜数C &  \\
{\sffamily\bfseries 4-9} & 波をたし合わせる：加法定理とうなり & \kgcstars{4}{4} & \kgcroute{4} & 数Ⅱ & \kgcout{p} \\
{\sffamily\bfseries 4-10} & 向きと大きさをもつ量：ベクトル & \kgcstars{2}{2} & \kgcroute{6} & 数C &  \\
{\sffamily\bfseries 4-11} & ベクトルを成分で計算しよう & \kgcstars{2}{3} & \kgcroute{6} & 数Ⅱ〜数C &  \\
{\sffamily\bfseries 4-12} & 内積：どれだけ同じ向きか & \kgcstars{3}{4} & \kgcroute{6} & 数Ⅰ〜数C &  \\
{\sffamily\bfseries 4-13} & 空間の座標と3次元のベクトル & \kgcstars{4}{4} & \kgcroute{6} & 小4〜数C &  \\
{\sffamily\bfseries 4-14} & 外積：回す力と右ねじ & \kgcstars{4}{5} & \kgcroute{6} & 大学 & \kgcout{a} \\
{\sffamily\bfseries 4-15} & いろいろな座標：極座標・円柱座標・球座標 & \kgcstars{4}{5} & \kgcroute{6} & 数C〜大学 & \kgcout{p} \\
\end{longtblr}
\Needspace{8\baselineskip}\kgbsub{第5巻}{変化をつかむ　微分と積分}\par\nobreak{\small\noindent 学年の目安：高校1年〜2年　章数：19}\par\nobreak\vspace{3pt}
\begin{longtblr}[theme=kgc, caption={第5巻　変化をつかむ}, entry=none, label=none]{colspec={Q[l,wd=14mm] X[l] Q[l,wd=17mm] Q[l,wd=12.5mm] Q[l,wd=25mm] Q[c,wd=10mm]}, width=\linewidth,
  rowhead=1, row{1}={bg=kgcHead, font=\sffamily\bfseries\footnotesize}, row{even[2]}={bg=kgcBand},
  hlines={0.4pt,bkLine}, vline{2-6}={0.4pt,bkLine}, rowsep=1.5pt, colsep=3pt, cells={font=\footnotesize, valign=m}, row{1}={valign=m}}
章 & 章名 & ★ & 路線 & 指導要領の学年・科目 & 範囲外 \\
{\sffamily\bfseries 5-0}\newline{\footnotesize\sffamily 序章} & スピードメーターは何を計算している？ & \kgcstars{2}{2} & \kgcroute{1} & ― &  \\
{\sffamily\bfseries 5-1} & 平均の速さと瞬間の速さ & \kgcstars{2}{2} & \kgcroute{1} & 小5〜数Ⅱ &  \\
{\sffamily\bfseries 5-2} & 時間をどんどん短くしてみよう & \kgcstars{2}{2} & \kgcroute{1} & 数Ⅱ &  \\
{\sffamily\bfseries 5-3} & かぎりなく近づく：極限と接線 & \kgcstars{3}{3} & \kgcroute{1} & 数Ⅱ〜数Ⅲ &  \\
{\sffamily\bfseries 5-4} & 表から導関数を見つけよう & \kgcstars{3}{3} & \kgcroute{1} & 数Ⅱ〜数Ⅲ &  \\
{\sffamily\bfseries 5-5} & 位置→速度→加速度 & \kgcstars{3}{3} & \kgcroute{1} & 数Ⅲ & \kgcout{p} \\
{\sffamily\bfseries 5-6} & 微分の道具箱：積と合成関数 & \kgcstars{4}{4} & \kgcroute{1} & 数Ⅲ &  \\
{\sffamily\bfseries 5-7} & いちばん高いのはどこ？　最大と最小 & \kgcstars{3}{3} & \kgcroute{1} & 数Ⅱ〜数Ⅲ &  \\
{\sffamily\bfseries 5-8} & 短い時間なら、かけ算で近似できる & \kgcstars{2}{2} & \kgcroute{1} & 範囲外 & \kgcout{a} \\
{\sffamily\bfseries 5-9} & グラフの面積は進んだ距離 & \kgcstars{2}{2} & \kgcroute{2} & 数Ⅱ &  \\
{\sffamily\bfseries 5-10} & 細い長方形に分けてたそう & \kgcstars{2}{4} & \kgcroute{2} & 数Ⅱ〜数Ⅲ &  \\
{\sffamily\bfseries 5-11} & 積分は微分の逆 & \kgcstars{3}{3} & \kgcroute{2} & 数Ⅱ〜数Ⅲ &  \\
{\sffamily\bfseries 5-12} & 物理での積分：仕事と流量 & \kgcstars{3}{3} & \kgcroute{2} & 数Ⅱ〜数Ⅲ &  \\
{\sffamily\bfseries 5-13} & 積分定数は「はじめの状態」 & \kgcstars{3}{3} & \kgcroute{2} & 数Ⅱ〜数Ⅲ & \kgcout{p} \\
{\sffamily\bfseries 5-14} & ふしぎな数 e & \kgcstars{3}{4} & \kgcroute{3} & 数Ⅲ &  \\
{\sffamily\bfseries 5-15} & sin・cosを微分すると？ & \kgcstars{4}{4} & \kgcroute{4} & 数Ⅲ & \kgcout{p} \\
{\sffamily\bfseries 5-16} & 積分の道具箱：置換積分と部分積分 & \kgcstars{4}{4} & \kgcroute{2} & 数Ⅲ &  \\
{\sffamily\bfseries 5-17} & 近くでは多項式で近似できる：テイラー展開 & \kgcstars{4}{4} & \kgcroute{1} & 大学 & \kgcout{a} \\
{\sffamily\bfseries 5-18} & まとめ：微分したときのふるまい & \kgcstars{3}{3} & \kgcroute{1} & 数Ⅱ〜大学 & \kgcout{p} \\
\end{longtblr}
\Needspace{8\baselineskip}\kgbsub{第6巻}{未来を計算する　微分方程式と複素数}\par\nobreak{\small\noindent 学年の目安：高校2年〜3年　章数：16}\par\nobreak\vspace{3pt}
\begin{longtblr}[theme=kgc, caption={第6巻　未来を計算する}, entry=none, label=none]{colspec={Q[l,wd=14mm] X[l] Q[l,wd=17mm] Q[l,wd=12.5mm] Q[l,wd=25mm] Q[c,wd=10mm]}, width=\linewidth,
  rowhead=1, row{1}={bg=kgcHead, font=\sffamily\bfseries\footnotesize}, row{even[2]}={bg=kgcBand},
  hlines={0.4pt,bkLine}, vline{2-6}={0.4pt,bkLine}, rowsep=1.5pt, colsep=3pt, cells={font=\footnotesize, valign=m}, row{1}={valign=m}}
章 & 章名 & ★ & 路線 & 指導要領の学年・科目 & 範囲外 \\
{\sffamily\bfseries 6-0}\newline{\footnotesize\sffamily 序章} & 「次の瞬間を決めるルール」で未来を予想する & \kgcstars{2}{2} & \kgcroute{1} & 大学 & \kgcout{a} \\
{\sffamily\bfseries 6-1} & 微分方程式ってなに？ & \kgcstars{2}{2} & \kgcroute{1} & 大学 & \kgcout{a} \\
{\sffamily\bfseries 6-2} & 表計算で解こう：オイラー法 & \kgcstars{2}{2} & \kgcroute{1} & 数B〜大学 & \kgcout{p} \\
{\sffamily\bfseries 6-3} & 積分するだけで解ける形 & \kgcstars{3}{3} & \kgcroute{2} & 数Ⅲ〜大学 & \kgcout{p} \\
{\sffamily\bfseries 6-4} & 自分に比例してへる：冷めるお湯 & \kgcstars{3}{3} & \kgcroute{3} & 大学 & \kgcout{a} \\
{\sffamily\bfseries 6-5} & 空気抵抗と終端速度 & \kgcstars{2}{3} & \kgcroute{1} & 大学 & \kgcout{a} \\
{\sffamily\bfseries 6-6} & ばねと振り子 & \kgcstars{3}{3} & \kgcroute{4} & 大学 & \kgcout{a} \\
{\sffamily\bfseries 6-7} & ゆれながら止まる：減衰振動 & \kgcstars{3}{4} & \kgcroute{4} & 大学 & \kgcout{a} \\
{\sffamily\bfseries 6-8} & 2乗してマイナスになる数：虚数 & \kgcstars{4}{4} & \kgcroute{4} & 数Ⅱ & \kgcout{p} \\
{\sffamily\bfseries 6-9} & 複素平面：かけ算は回転と拡大 & \kgcstars{4}{4} & \kgcroute{4} & 数C &  \\
{\sffamily\bfseries 6-10} & オイラーの公式 & \kgcstars{4}{5} & \kgcroute{4} & 大学 & \kgcout{a} \\
{\sffamily\bfseries 6-11} & 双曲線関数 cosh・sinh & \kgcstars{4}{5} & \kgcroute{3} & 大学 & \kgcout{a} \\
{\sffamily\bfseries 6-12} & 減衰振動を式で解く：特性方程式 & \kgcstars{5}{5} & \kgcroute{4} & 大学 & \kgcout{a} \\
{\sffamily\bfseries 6-13} & ゆらし続けると：強制振動と共振 & \kgcstars{5}{5} & \kgcroute{4} & 大学 & \kgcout{a} \\
{\sffamily\bfseries 6-14} & 微分方程式がいくつも：連立微分方程式 & \kgcstars{5}{5} & \kgcroute{5} & 大学 & \kgcout{a} \\
{\sffamily\bfseries 6-15} & まとめ：物理は「微分方程式＋はじめの状態」で書ける & \kgcstars{2}{3} & \kgcroute{1} & 大学 & \kgcout{a} \\
\end{longtblr}
\Needspace{8\baselineskip}\kgbsub{第7巻}{まとめて計算する　行列}\par\nobreak{\small\noindent 学年の目安：高校3年（発展）〜大学1年　章数：12}\par\nobreak\vspace{3pt}
\begin{longtblr}[theme=kgc, caption={第7巻　まとめて計算する}, entry=none, label=none]{colspec={Q[l,wd=14mm] X[l] Q[l,wd=17mm] Q[l,wd=12.5mm] Q[l,wd=25mm] Q[c,wd=10mm]}, width=\linewidth,
  rowhead=1, row{1}={bg=kgcHead, font=\sffamily\bfseries\footnotesize}, row{even[2]}={bg=kgcBand},
  hlines={0.4pt,bkLine}, vline{2-6}={0.4pt,bkLine}, rowsep=1.5pt, colsep=3pt, cells={font=\footnotesize, valign=m}, row{1}={valign=m}}
章 & 章名 & ★ & 路線 & 指導要領の学年・科目 & 範囲外 \\
{\sffamily\bfseries 7-0}\newline{\footnotesize\sffamily 序章} & 表のままで計算できたら？ & \kgcstars{3}{3} & \kgcroute{5} & 数C〜大学 & \kgcout{p} \\
{\sffamily\bfseries 7-1} & 数の表で計算する：行列 & \kgcstars{3}{3} & \kgcroute{5} & 数C〜大学 & \kgcout{p} \\
{\sffamily\bfseries 7-2} & 行列×ベクトル：形を変える機械 & \kgcstars{4}{4} & \kgcroute{5} & 大学 & \kgcout{a} \\
{\sffamily\bfseries 7-3} & 行列×行列：変換をつなぐ & \kgcstars{4}{4} & \kgcroute{5} & 大学 & \kgcout{a} \\
{\sffamily\bfseries 7-4} & 行列式：面積が何倍になるか & \kgcstars{4}{4} & \kgcroute{5} & 大学 & \kgcout{a} \\
{\sffamily\bfseries 7-5} & 逆行列で連立方程式を解く & \kgcstars{4}{4} & \kgcroute{5} & 大学 & \kgcout{a} \\
{\sffamily\bfseries 7-6} & 回転行列：くるっと回す & \kgcstars{4}{4} & \kgcroute{4} & 数Ⅱ〜大学 & \kgcout{p} \\
{\sffamily\bfseries 7-7} & 3次元の回転と座標の変換 & \kgcstars{5}{5} & \kgcroute{6} & 大学 & \kgcout{a} \\
{\sffamily\bfseries 7-8} & 固有値と固有ベクトル：向きの変わらない方向 & \kgcstars{5}{5} & \kgcroute{5} & 大学 & \kgcout{a} \\
{\sffamily\bfseries 7-9} & 対角化：まっすぐな方向に分けて考える & \kgcstars{5}{5} & \kgcroute{5} & 大学 & \kgcout{a} \\
{\sffamily\bfseries 7-10} & 行列のくり返し：漸化式とオイラー法を行列で & \kgcstars{5}{5} & \kgcroute{5} & 数B〜大学 & \kgcout{p} \\
{\sffamily\bfseries 7-11}\newline{\footnotesize\sffamily 読みもの} & 行列の使いみち & \kgcstars{3}{5} & \kgcroute{5} & 大学 & \kgcout{a} \\
\end{longtblr}
\Needspace{8\baselineskip}\kgbsub{第8巻}{空間に広がる量　ベクトル解析}\par\nobreak{\small\noindent 学年の目安：大学1年〜2年　章数：15}\par\nobreak\vspace{3pt}
\begin{longtblr}[theme=kgc, caption={第8巻　空間に広がる量}, entry=none, label=none]{colspec={Q[l,wd=14mm] X[l] Q[l,wd=17mm] Q[l,wd=12.5mm] Q[l,wd=25mm] Q[c,wd=10mm]}, width=\linewidth,
  rowhead=1, row{1}={bg=kgcHead, font=\sffamily\bfseries\footnotesize}, row{even[2]}={bg=kgcBand},
  hlines={0.4pt,bkLine}, vline{2-6}={0.4pt,bkLine}, rowsep=1.5pt, colsep=3pt, cells={font=\footnotesize, valign=m}, row{1}={valign=m}}
章 & 章名 & ★ & 路線 & 指導要領の学年・科目 & 範囲外 \\
{\sffamily\bfseries 8-0}\newline{\footnotesize\sffamily 序章} & 空間のどこにも矢印がある & \kgcstars{3}{3} & \kgcroute{6} & 大学 & \kgcout{a} \\
{\sffamily\bfseries 8-1} & 地図と等高線：2つ以上の変数の関数 & \kgcstars{4}{4} & \kgcroute{1} & 大学 & \kgcout{a} \\
{\sffamily\bfseries 8-2} & 偏微分：ほかを止めて1つだけ動かす & \kgcstars{4}{4} & \kgcroute{1} & 大学 & \kgcout{a} \\
{\sffamily\bfseries 8-3} & 全微分：少し動いたときの変化 & \kgcstars{5}{5} & \kgcroute{1} & 大学 & \kgcout{a} \\
{\sffamily\bfseries 8-4} & いちばん急な向き：勾配 grad & \kgcstars{5}{5} & \kgcroute{1} & 大学 & \kgcout{a} \\
{\sffamily\bfseries 8-5} & ベクトル場 & \kgcstars{4}{4} & \kgcroute{6} & 大学 & \kgcout{a} \\
{\sffamily\bfseries 8-6} & 面や立体でたし合わせる：重積分 & \kgcstars{5}{5} & \kgcroute{2} & 大学 & \kgcout{a} \\
{\sffamily\bfseries 8-7} & 道にそってたす：線積分 & \kgcstars{5}{5} & \kgcroute{2} & 数Ⅲ〜大学 & \kgcout{p} \\
{\sffamily\bfseries 8-8} & 面を通りぬける量：面積分と流束 & \kgcstars{5}{5} & \kgcroute{2} & 大学 & \kgcout{a} \\
{\sffamily\bfseries 8-9} & わき出しと吸いこみ：発散 div & \kgcstars{5}{5} & \kgcroute{6} & 大学 & \kgcout{a} \\
{\sffamily\bfseries 8-10} & うず：回転 rot & \kgcstars{5}{5} & \kgcroute{6} & 大学 & \kgcout{a} \\
{\sffamily\bfseries 8-11} & 中と外をつなぐ：ガウスの定理とストークスの定理 & \kgcstars{5}{5} & \kgcroute{2} & 大学 & \kgcout{a} \\
{\sffamily\bfseries 8-12} & ならす働き：ラプラシアン & \kgcstars{5}{5} & \kgcroute{1} & 大学 & \kgcout{a} \\
{\sffamily\bfseries 8-13} & 熱と波の伝わり方：偏微分方程式 & \kgcstars{5}{5} & \kgcroute{1} & 大学 & \kgcout{a} \\
{\sffamily\bfseries 8-14}\newline{\footnotesize\sffamily 読みもの} & div と rot で書かれた式を読もう & \kgcstars{5}{5} & \kgcroute{6} & 大学 & \kgcout{a} \\
\end{longtblr}
\Needspace{8\baselineskip}\kgbsub{第9巻}{波と変化を読みとく　フーリエ解析とラプラス変換}\par\nobreak{\small\noindent 学年の目安：大学2年　章数：14}\par\nobreak\vspace{3pt}
\begin{longtblr}[theme=kgc, caption={第9巻　波と変化を読みとく}, entry=none, label=none]{colspec={Q[l,wd=14mm] X[l] Q[l,wd=17mm] Q[l,wd=12.5mm] Q[l,wd=25mm] Q[c,wd=10mm]}, width=\linewidth,
  rowhead=1, row{1}={bg=kgcHead, font=\sffamily\bfseries\footnotesize}, row{even[2]}={bg=kgcBand},
  hlines={0.4pt,bkLine}, vline{2-6}={0.4pt,bkLine}, rowsep=1.5pt, colsep=3pt, cells={font=\footnotesize, valign=m}, row{1}={valign=m}}
章 & 章名 & ★ & 路線 & 指導要領の学年・科目 & 範囲外 \\
{\sffamily\bfseries 9-0}\newline{\footnotesize\sffamily 序章} & 音を分ける、変化を予想する & \kgcstars{3}{3} & \kgcroute{4} & 大学 & \kgcout{a} \\
{\sffamily\bfseries 9-1} & 波をたし合わせて形を作る：フーリエ級数 & \kgcstars{4}{4} & \kgcroute{4} & 大学 & \kgcout{a} \\
{\sffamily\bfseries 9-2} & 成分の取り出し方：直交性 & \kgcstars{5}{5} & \kgcroute{4} & 大学 & \kgcout{a} \\
{\sffamily\bfseries 9-3} & 複素数で書くフーリエ級数 & \kgcstars{5}{5} & \kgcroute{4} & 大学 & \kgcout{a} \\
{\sffamily\bfseries 9-4} & スペクトル：フーリエ変換 & \kgcstars{5}{5} & \kgcroute{7} & 大学 & \kgcout{a} \\
{\sffamily\bfseries 9-5} & 微分をかけ算に変える翻訳機：ラプラス変換 & \kgcstars{5}{5} & \kgcroute{3} & 大学 & \kgcout{a} \\
{\sffamily\bfseries 9-6} & 変換表と性質 & \kgcstars{5}{5} & \kgcroute{3} & 大学 & \kgcout{a} \\
{\sffamily\bfseries 9-7} & ラプラス変換で微分方程式を解く & \kgcstars{5}{5} & \kgcroute{3} & 数Ⅱ〜大学 & \kgcout{p} \\
{\sffamily\bfseries 9-8} & 伝達関数：入れると出てくるを s で表す & \kgcstars{5}{5} & \kgcroute{3} & 大学 & \kgcout{a} \\
{\sffamily\bfseries 9-9} & 極と安定性：s の値でふるまいがわかる & \kgcstars{5}{5} & \kgcroute{3} & 数Ⅱ〜大学 & \kgcout{p} \\
{\sffamily\bfseries 9-10} & 連立微分方程式を行列で：状態方程式 & \kgcstars{5}{5} & \kgcroute{5} & 大学 & \kgcout{a} \\
{\sffamily\bfseries 9-11} & 行列の指数関数と固有値 & \kgcstars{5}{5} & \kgcroute{5} & 数Ⅲ〜大学 & \kgcout{p} \\
{\sffamily\bfseries 9-12}\newline{\footnotesize\sffamily 読みもの} & フーリエとラプラスの使いみち & \kgcstars{3}{5} & \kgcroute{7} & 大学 & \kgcout{a} \\
{\sffamily\bfseries 9-13}\newline{\footnotesize\sffamily 終章} & 工学の数学の地図 & \kgcstars{3}{5} & \kgcallroutes{} & ― &  \\
\end{longtblr}
